% paper.tex
% ============================================================================
% Unified working paper for the project on two graphon optimisation problems.
%
% This single file consolidates the entire current state of the project.  It
% merges, with all mathematics preserved verbatim:
%   * two_problems_progress_report.tex  (the overall progress report)
%   * Tk_first_order_stability.tex      (rigorous first-order stability of T_k)
%   * rigorous_certificates.tex         (Sturm certificates for Psi_P, Psi_J)
%   * balancing_reduction.tex           (Step A2: in-family balancing, corrected)
%   * global_reduction_strategy.tex     (research roadmap; Plans A/B/C)
% and additionally records the direct step-graphon optimisation experiment
% (scripts/optimize_graphon.py) that provides numerical evidence for Step A1.
%
% Every algebraic certificate is reproduced by the companion scripts in
% scripts/ (run with /opt/miniconda3/bin/python3, which has sympy 1.14);
% see Appendix C for the script index.
%
% The individual source .tex files are kept for now but this file is intended
% to be the single point of reference for collaborators.
% ============================================================================

\documentclass[11pt]{amsart}

\usepackage{amsmath,amssymb,amsthm,mathtools}
\usepackage[margin=1.1in]{geometry}
\usepackage{enumitem}
\usepackage{hyperref}
\usepackage{xcolor}

\hypersetup{
  colorlinks=true,
  linkcolor=blue!50!black,
  citecolor=blue!50!black,
  urlcolor=blue!50!black
}

\newtheorem{theorem}{Theorem}[section]
\newtheorem{proposition}[theorem]{Proposition}
\newtheorem{lemma}[theorem]{Lemma}
\newtheorem{conjecture}[theorem]{Conjecture}
\newtheorem{problem}[theorem]{Problem}
\newtheorem{corollary}[theorem]{Corollary}

\theoremstyle{definition}
\newtheorem{definition}[theorem]{Definition}
\newtheorem{remark}[theorem]{Remark}
\newtheorem{example}[theorem]{Example}
\newtheorem{status}[theorem]{Status}

\newcommand{\graphon}{\mathcal W}
\newcommand{\E}{\mathbb E}
\newcommand{\1}{\mathbf 1}
\newcommand{\R}{\mathbb R}
\newcommand{\Z}{\mathbb Z}
\newcommand{\dd}{\,d}
\newcommand{\tk}{T_k}
\newcommand{\LS}{\mathrm{LS}}
\newcommand{\dist}{\operatorname{dist}}
\newcommand{\Hom}{\operatorname{Hom}}
\newcommand{\ord}{\operatorname{ord}}
\DeclareMathOperator{\Hess}{Hess}

\title[Two graphon optimisation problems: a consolidated progress paper]%
{Two graphon optimisation problems:\\
counterexamples, corrected conjectures, reduced variational targets,\\
and rigorous certificates}
\author{Working research paper (consolidated)}
\date{\today}

\begin{document}

\begin{abstract}
This paper consolidates the current state of two graphon optimisation problems.
The first asks whether, for a particular six-vertex graph $H$, the minimiser of
$t(H,W)$ at edge density $1-1/k$ is never the balanced $k$-partite graphon nor
the constant graphon, and asks for an exact solution at $k=5$.  We verify the
balanced and constant benchmark values, prove the proposed $k=5$ circular
construction has density $24349/187500$, and give explicit competitors beating
both benchmark graphons for $k=3,4,5$.  We also prove --- with the averaging
argument made fully explicit --- that for $k\ge 6$ the balanced $k$-partite
graphon is first-order locally stable against all edge-density-preserving
feasible perturbations, so the all-$k$ statement is not accessible by a local
perturbation.

The second problem asks for a convex/concave classification of the extremal
curves $f_F(x)=\min\{t(F,W):t(K_2,W)=x\}$.  We show that the proposed extension
to all larger graphs is false: $F=K_3\sqcup K_3$ already gives a curve that is
neither convex nor concave on $[1/2,2/3]$, and we formulate a repaired
``odd-prime'' direction.  The clique subfamily is solved by Reiher's
clique-density theorem and is piecewise concave.  However, two non-clique
cases listed as concave in the original problem, namely $P=K_4^\dagger$ ($K_4$
with a pendant edge) and $J=K_3\cup_{K_2}K_4$, do not have the
Lov\'asz--Simonovits multipartite graphon as their minimiser: we exhibit
explicit better tripartite Tur\'an-filling graphons.  We derive the exact
finite-dimensional reduced variational curves $\Psi_P,\Psi_J$ for the filling
family, and prove a battery of \emph{rigorous Sturm-sequence certificates}:
branch uniqueness, monotonicity of the parametrisation, endpoint values, the
strict beating of the feasibility boundary, and --- the substantive new finding
--- that $\Psi_P$ and $\Psi_J$ each have \emph{mixed convex-then-concave
curvature} on the first Tur\'an interval $[2/3,3/4]$, with unique inflection
points $x^*_P\approx 0.6788$ and $x^*_J\approx 0.7132$.

We then close \emph{Step A2} of the global reduction --- the reduction of the
unbalanced tripartite-filling family to balanced parts --- in its corrected
global form (the natural per-$\alpha$ balancing statement is false, and we give
an explicit counterexample).  We record the research roadmap for the remaining
\emph{Step A1} (cut-metric stability forcing a $3$-blowup skeleton), and present
a direct step-graphon optimisation experiment whose output supports both the
global-min conjecture $f_P=\Psi_P$ and the Step~A1 stability mechanism.

All algebraic certificates are reproduced by the companion \texttt{sympy}
scripts indexed in Appendix~\ref{app:scripts}.
\end{abstract}

\maketitle

\tableofcontents

% ============================================================================
\section{Introduction and reading guide}
\label{sec:intro}

This is a single-file consolidation of the project's progress notes.  It is
intended as the one document a collaborator needs to read to know what has been
achieved.  The logical structure is:

\begin{itemize}[itemsep=3pt]
\item \textbf{Problem 1} (the six-vertex graph $H$): \S\ref{sec:problem1} records
the benchmark values, the circular $k=5$ construction, and explicit competitors;
\S\ref{sec:tkstability} proves first-order stability of $T_k$ for $k\ge 6$.
\item \textbf{Problem 2} (the convex/concave classification):
\S\ref{sec:problem2false} disproves the all-graphs extension;
\S\ref{sec:oddprime} proposes the odd-prime repair;
\S\ref{sec:clique} records the solved clique subfamily;
\S\ref{sec:nonclique} exhibits the two non-clique counterexamples
$P=K_4^\dagger$ and $J=K_3\cup_{K_2}K_4$;
\S\ref{sec:filling} derives the reduced curves $\Psi_P,\Psi_J$;
\S\ref{sec:certificates} proves the rigorous Sturm certificates for them,
including the mixed-curvature theorem;
\S\ref{sec:global} lays out the global-reduction roadmap;
\S\ref{sec:balancing} closes Step~A2;
\S\ref{sec:numerics} reports the step-graphon optimisation experiment.
\item \textbf{Status and directions}: \S\ref{sec:status}--\S\ref{sec:checklist}.
\end{itemize}

\paragraph{What is rigorous.}  The following are proved here (and machine-checked
where indicated): the benchmark and competitor values for Problem~1
(\S\ref{sec:problem1}); first-order stability of $T_k$ for $k\ge6$
(Theorem~\ref{thm:first-order-stability}); the $K_3\sqcup K_3$ counterexample
(\S\ref{sec:problem2false}); the clique-curve concavity
(Theorem~\ref{thm:clique-concave}); strict improvement over the LS template at
$x=17/25$ (\S\ref{sec:nonclique}); the in-family reduction to $\Psi_P,\Psi_J$
(\S\ref{sec:filling}); branch uniqueness, monotonicity, endpoints, mixed
curvature, and the stationary-vs-boundary gap
(Theorems~\ref{thm:branch-range}--\ref{thm:stat-vs-bdry}); and the corrected
in-family balancing reduction, i.e.\ Step~A2
(Theorem~\ref{thm:balancing-main}), modulo a single Sturm endgame flagged in
\S\ref{sec:balancing-status}.  The remaining genuinely open target for
$f_P=\Psi_P$ globally is \emph{Step~A1}; \S\ref{sec:numerics} gives numerical
evidence for it.

% ============================================================================
\section{Notation and basic conventions}
\label{sec:notation}

A \emph{graphon} is a symmetric measurable function
\[
   W:[0,1]^2\to[0,1].
\]
For a finite graph $F=(V(F),E(F))$, its homomorphism density in $W$ is
\[
   t(F,W)=\int_{[0,1]^{V(F)}}\prod_{ij\in E(F)} W(x_i,x_j)\prod_{i\in V(F)}\dd x_i.
\]
In particular,
\[
   t(K_2,W)=\int_{[0,1]^2}W(x,y)\dd x\dd y
\]
is the ordered-pair edge density.  For a fixed graph $F$, define
\[
   f_F(x)=\min\{t(F,W):t(K_2,W)=x\},\qquad x\in[0,1].
\]
Existence of a minimiser follows from standard graphon compactness and lower
semicontinuity of homomorphism densities.

The balanced complete $k$-partite graphon will be denoted by $T_k$.  It is
obtained by partitioning $[0,1]$ into $k$ equal parts and putting value $0$ on
each diagonal part and value $1$ between distinct parts.  Thus
\[
   t(K_2,T_k)=1-\frac1k.
\]
The constant graphon of edge density $p$ is denoted by $Q_p$; thus
$Q_p(x,y)=p$ and
\[
   t(F,Q_p)=p^{e(F)}.
\]

We write
\[
   I_k=\left[1-\frac1k,\;1-\frac1{k+1}\right]
\]
for the usual Tur\'an intervals.

% ============================================================================
\section{Problem 1: the six-vertex graph}
\label{sec:problem1}

Let $H$ be the graph with vertex set $\{1,2,3,4,5,6\}$ and edge set
\[
\begin{split}
E(H)=\{&13,14,23,24,35,36,45,46,56\}.
\end{split}
\]
Thus $H$ has six vertices and nine edges.  The first problem asks about
\[
   \min\{t(H,W):t(K_2,W)=1-1/k\}.
\]

\subsection{The balanced $k$-partite and constant values}

\begin{lemma}[Balanced $k$-partite value]
For every integer $k\ge 1$,
\[
   t(H,T_k)=\frac{k(k-1)(k-2)(k^3-6k^2+14k-11)}{k^6}.
\]
\end{lemma}

\begin{proof}
A homomorphism from $H$ to $T_k$ is exactly a proper $k$-colouring of $H$.  First
choose colours of vertices $5$ and $6$.  They must be distinct, giving $k(k-1)$
choices.  The vertices $3$ and $4$ must avoid both colours used on $5$ and $6$.

There are two cases.  If $3$ and $4$ receive the same colour, there are $k-2$
choices for this colour; then vertices $1$ and $2$ may each use any colour
except that colour, giving $(k-1)^2$ choices.  If $3$ and $4$ receive distinct
colours, there are $(k-2)(k-3)$ choices; then vertices $1$ and $2$ may each use
any colour except the two colours used on $3,4$, giving $(k-2)^2$ choices.
Hence
\[
\begin{split}
\chi_H(k)&=k(k-1)\bigl((k-2)(k-1)^2+(k-2)(k-3)(k-2)^2\bigr)\\
&=k(k-1)(k-2)(k^3-6k^2+14k-11).
\end{split}
\]
Dividing by $k^6$ gives the stated homomorphism density.
\end{proof}

\begin{lemma}[Constant value]
For $Q_{1-1/k}\equiv 1-1/k$,
\[
   t(H,Q_{1-1/k})=\left(1-\frac1k\right)^9.
\]
\end{lemma}

\begin{proof}
The graph $H$ has nine edges, and every edge contributes the constant factor
$1-1/k$.
\end{proof}

For reference, the first few benchmark values are
\[
\begin{array}{c|c|c}
 k & t(H,T_k) & t(H,Q_{1-1/k})\\ \hline
 3 & 8/243\approx 0.0329218 & (2/3)^9\approx 0.0260123\\
 4 & 39/512\approx 0.0761719 & (3/4)^9\approx 0.0750847\\
 5 & 2040/15625=0.13056 & (4/5)^9\approx 0.1342177.
\end{array}
\]

\subsection{The proposed circular construction at $k=5$}

Define
\[
   W_\circ(x,y)=\1\{0.1\le |x-y|\le 0.9\}.
\]
Equivalently, the zero set is the circular interval of radius $1/10$ around the
diagonal.  Its measure is $1/5$, so
\[
   t(K_2,W_\circ)=\frac45.
\]

\begin{theorem}[Density of the circular construction]
The circular graphon $W_\circ$ satisfies
\[
   t(H,W_\circ)=\frac{24349}{187500}\approx 0.129861333\ldots .
\]
Consequently
\[
   t(H,W_\circ)<t(H,T_5)<t(H,Q_{4/5}).
\]
\end{theorem}

\begin{proof}
We approximate $W_\circ$ by finite cyclic step graphons.  Let $d$ be odd and set
$N=5d$.  Partition $[0,1]$ into $N$ equal intervals indexed by $\Z/N\Z$.  Let
$W_d$ be the step graphon with matrix
\[
   W_d(i,j)=1\quad\Longleftrightarrow\quad
   \dist_{C_N}(i,j)>\frac{d-1}{2}.
\]
For every row there are exactly $d$ zero entries, including the diagonal entry.
Hence
\[
   t(K_2,W_d)=1-\frac{d}{5d}=\frac45.
\]
Moreover $W_d\to W_\circ$ in $L^1$, and therefore $t(H,W_d)\to t(H,W_\circ)$.

A direct cyclic count gives
\[
   \#\Hom(H,W_d)=\frac{24349d^6+140d^4-9d^2}{12}.
\]
Here $\Hom(H,W_d)$ means the number of maps $\phi:V(H)\to\Z/N\Z$ such that all
nine required pairs are $W_d$-edges.  Dividing by $N^6=(5d)^6$ gives
\[
   t(H,W_d)=\frac{24349+140d^{-2}-9d^{-4}}{187500}.
\]
Letting $d\to\infty$ gives
\[
   t(H,W_\circ)=\frac{24349}{187500}.
\]
Finally,
\[
   \frac{24349}{187500}<\frac{2040}{15625}<\left(\frac45\right)^9.
\]
Indeed $2040/15625=24480/187500$, and $24480>24349$.
\end{proof}

\begin{remark}[What this proves and what it does not prove]
The theorem proves that the balanced $5$-partite and constant graphons are not
optimal at edge density $4/5$.  It does \emph{not} prove that $W_\circ$ is
globally optimal.  The missing statement is the lower bound
\[
   t(H,W)\ge \frac{24349}{187500}\qquad
   \text{whenever }t(K_2,W)=4/5.
\]
This lower bound is the hard part of the $k=5$ problem and is a natural
flag-algebra target.
\end{remark}

\subsection{Explicit competitors for $k=3,4,5$}
\label{ssec:Ska}

There is a simple one-parameter family that beats both the balanced and constant
graphons for $k=3,4,5$.

Partition $[0,1]$ into $k$ equal parts and define
\[
   S_{k,a}(x,y)=
   \begin{cases}
      a, & x,y\text{ in the same part},\\[1mm]
      1-\dfrac{a}{k-1}, & x,y\text{ in different parts}.
   \end{cases}
\]
Then
\[
   t(K_2,S_{k,a})=1-\frac1k
\]
for every $a\in[0,1]$ for which $1-a/(k-1)\in[0,1]$.

For this family,
\[
   t(H,S_{k,a})=\frac1{k^6}\sum_{c_1,\ldots,c_6\in[k]}
        \prod_{ij\in E(H)}w(c_i,c_j),
\]
where
\[
   w(c_i,c_j)=
   \begin{cases}
      a, & c_i=c_j,\\
      1-a/(k-1), & c_i\ne c_j.
   \end{cases}
\]
Exact rational evaluation gives:
\[
\begin{array}{c|c|c|c}
 k & a & t(H,S_{k,a}) & \text{comparison}\\ \hline
 3 & 3/10 & \dfrac{792167976529}{31104000000000}\approx 0.02546836
   & <(2/3)^9<t(H,T_3)\\[2mm]
 4 & 1/5 & \dfrac{38420325073}{524880000000}\approx 0.07319830
   & <(3/4)^9<t(H,T_4)\\[2mm]
 5 & 1/20 & \dfrac{2186447178930623}{16777216000000000}\approx 0.13032241
   & <t(H,T_5)<(4/5)^9.
\end{array}
\]
Thus the statement ``neither balanced nor constant'' is rigorously verified for
$k=3,4,5$.

\subsection{Why the all-$k$ statement needs a global argument}

The one-parameter perturbation above also explains why the problem changes at
$k=6$.

\begin{proposition}[First-order local stability of $T_k$ for $k\ge 6$]
\label{prop:tk-stability-short}
For the family $S_{k,a}$,
\[
   \left.\frac{d}{da}t(H,S_{k,a})\right|_{a=0}
   =\frac{6k^3-55k^2+142k-111}{k^5}.
\]
This derivative is negative for $k=3,4,5$ and positive for every $k\ge 6$.
Consequently, for $k\ge 6$, the balanced $k$-partite graphon is first-order
locally stable against every edge-density-preserving feasible perturbation.
\end{proposition}

\begin{proof}
The displayed derivative follows by differentiating the finite sum for
$t(H,S_{k,a})$ at $a=0$.  The polynomial $6k^3-55k^2+142k-111$ is negative at
$k=3,4,5$ and positive at $k=6$; its derivative $18k^2-110k+142$ is positive for
$k\ge 5$, so the polynomial remains positive for all $k\ge 6$.

That this one-parameter derivative controls \emph{all} first-order feasible
directions at $T_k$ is exactly the averaging statement proved in full in
\S\ref{sec:tkstability} (Theorem~\ref{thm:first-order-stability}): by the
symmetry of $T_k$, every edge-density-preserving feasible perturbation can be
$G_k$-symmetrised without changing the first variation, and the symmetric
feasibility cone is the single half-line tangent to $\{S_{k,a}\}_{a\ge0}$.
\end{proof}

\begin{status}[Status of Problem 1]
The following points are solid.
\begin{enumerate}[label=(\roman*)]
\item The balanced and constant benchmark values are explicit.
\item At $k=5$, the circular graphon has the proposed value $24349/187500$ and
beats both benchmarks.
\item For $k=3,4,5$, explicit rational competitors beat both benchmarks.
\item For $k\ge 6$, non-optimality of $T_k$ cannot be proved by an infinitesimal
edge-density-preserving perturbation, because $T_k$ is first-order locally
stable (Theorem~\ref{thm:first-order-stability}).
\end{enumerate}
The missing results are global lower and upper-bound statements.  The two most
important targets are:
\begin{enumerate}[label=(\alph*)]
\item prove or disprove global optimality of the circular $k=5$ graphon;
\item find a genuinely global competitor, or a proof of optimality, for
$k\ge 6$.
\end{enumerate}
\end{status}

\begin{remark}[Normalisation of circular thresholds]
To obtain edge density $1-1/k$ from a circular interval construction, the
forbidden circular radius should be $1/(2k)$, not $1/k$.  Thus the literal
analogue of the $k=5$ thresholds $0.1,0.9$ would use thresholds
\[
   \frac1{2k},\qquad 1-\frac1{2k}.
\]
For $k=3,4$, these would be $1/6,5/6$ and $1/8,7/8$, respectively.  This explains
why replacing $0.1,0.9$ by $1/3,2/3$ or $1/4,3/4$ does not preserve the edge
density.
\end{remark}

% ============================================================================
\section{First-order stability of $T_k$ for $k\ge 6$: the averaging step}
\label{sec:tkstability}

This section proves Proposition~\ref{prop:tk-stability-short} in full, with the
averaging argument that reduces all feasible directions to the family
$\{S_{k,a}\}$ made completely explicit.  The statement is that, for the
six-vertex graph $H$ and every integer $k\ge 6$, the balanced complete
$k$-partite graphon $T_k$ is first-order stable under the edge-density
constraint $t(K_2,W)=1-1/k$: no first-order feasible perturbation of $T_k$
decreases $t(H,W)$.

\subsection{Admissible first-order perturbations at $T_k$}

With the partition $[0,1] = I_1 \sqcup \cdots \sqcup I_k$ into $k$ equal
intervals $I_j = [(j-1)/k, j/k]$, we have $T_k(x,y) = 1$ iff $x,y$ lie in
different parts, and $T_k(x,y) = 0$ iff $x,y$ lie in the same part.

For $f\in L^\infty([0,1]^2)$ symmetric, write $W_t := T_k + t\, f$.  The
\emph{first-order feasibility cone at $T_k$ under the edge-density constraint}
$t(K_2,W) = 1-1/k$ is the closed convex cone
\begin{equation}
\label{eq:feas-cone}
\mathcal F_k := \left\{
   f : [0,1]^2 \to \R\ \text{measurable, symmetric, } \|f\|_\infty<\infty
   \ \middle|\
   \begin{aligned}
   &f(x,y)\ge 0\ \text{a.e.\ on within-part cells},\\
   &f(x,y)\le 0\ \text{a.e.\ on cross-part cells},\\
   &\int_{[0,1]^2} f \dd x\dd y = 0
   \end{aligned}
\right\}.
\end{equation}
The first two conditions ensure that for $t>0$ small enough, $W_t$ takes values
in $[0,1]$ (so $W_t$ is a graphon), and the third ensures that
$t(K_2,W_t)=1-1/k+0\cdot t + O(t^2)$.

\begin{remark}
Strictly, every ``feasible perturbation'' should be defined modulo the
equivalence of graphons under measure-preserving transformations.  All
expressions below depend only on the equivalence class, so we may work with
representatives.
\end{remark}

\subsection{The first variation kernel and its symmetry}

For a finite graph $H$ and a graphon $W$, the standard expansion of $t(H,W)$
gives the first variation at $W=T_k$ in the direction $f$:
\begin{equation}
\label{eq:first-variation}
   \delta t(H,T_k)[f]
   = \sum_{ij\in E(H)}
     \int_{[0,1]^2} \mathcal K_{ij}(x,y)\, f(x,y) \dd x\dd y,
\end{equation}
where
\[
   \mathcal K_{ij}(x,y)
   = \int_{[0,1]^{V(H)\setminus\{i,j\}}}
       \prod_{\substack{i'j'\in E(H) \\ \{i',j'\}\ne\{i,j\}}} T_k(x_{i'},x_{j'})
       \prod_{v\in V(H)\setminus\{i,j\}} \dd x_v,
\]
with $x_i:=x$ and $x_j:=y$.  Symmetrising over $E(H)$, define the total kernel
$\mathcal K(x,y) := \sum_{ij\in E(H)} \mathcal K_{ij}(x,y)$ (symmetric in
$x,y$).  Then
\[
   \delta t(H,T_k)[f] = \int_{[0,1]^2} \mathcal K(x,y) f(x,y) \dd x\dd y.
\]

\begin{lemma}[Block-constancy of the first-variation kernel]
\label{lem:K-block-constant}
There exist constants $c_w, c_c\in\R$ depending only on $H$ and $k$ such that
\[
   \mathcal K(x,y) =
   \begin{cases}
      c_w, & x,y\ \text{in the same part of }T_k,\\
      c_c, & x,y\ \text{in different parts of }T_k,
   \end{cases}
   \qquad\text{a.e.\ on }[0,1]^2.
\]
\end{lemma}

\begin{proof}
The kernel $\mathcal K(x,y)$ depends on $(x,y)$ only through the product
$\prod T_k(x_{i'}, x_{j'})$ for various tuples, which is an indicator of a
proper $k$-colouring condition.  Fix $(x,y)$ in the support of $\mathcal K$ and
let $j(x), j(y)\in\{1,\ldots,k\}$ be the parts containing $x$ and $y$
respectively.  For any block permutation $\sigma\in S_k$ that fixes $j(x)$ and
$j(y)$, the measure-preserving map of $[0,1]$ induced by $\sigma$ preserves
$T_k$ and hence preserves $\mathcal K$.  Combined with the $S_k$-action, this
reduces the kernel value at $(x,y)$ to a function of only the pair
$\{j(x), j(y)\}$, and indeed only to the equivalence class ``same part'' or
``different parts'' (since any two same-part pairs are related by a block
permutation, and any two different-part pairs are related by one).
\end{proof}

\subsection{The averaging step}

The space $\mathcal F_k$ in~\eqref{eq:feas-cone} is naturally acted upon by the
group $G_k$ consisting of pairs $(\sigma, \tau)$ where $\sigma\in S_k$ is a
permutation of the parts and $\tau$ is a measure-preserving automorphism of
$[0,1]$ preserving the partition $\{I_1,\ldots,I_k\}$.  Both $T_k$ and the
edge-density constraint are $G_k$-invariant, so the action preserves
$\mathcal F_k$ and preserves $\delta t(H,T_k)[\cdot]$.

\begin{definition}[Symmetrised perturbation]
For $f\in\mathcal F_k$, let
\[
   \bar f(x,y) := \frac{1}{|S_k|}\sum_{\sigma\in S_k}
   \int_{\mathrm{Aut}_\partial}\ (f^{(\sigma,\tau)})(x,y)\dd\mu(\tau),
\]
where $\mathrm{Aut}_\partial$ denotes the group of partition-preserving
measure-preserving transformations of $[0,1]$ and $\mu$ is its Haar measure
(equivalently, the product over blocks of the Haar measure on the
measure-preserving group of each $I_j$).  Equivalently, $\bar f$ is the
projection of $f$ onto the $G_k$-invariant subspace of $L^2([0,1]^2)$.
\end{definition}

\begin{lemma}[Averaging preserves feasibility and the linear functional]
\label{lem:averaging}
For every $f\in\mathcal F_k$, the symmetrised perturbation $\bar f$ also lies in
$\mathcal F_k$, and
\[
   \delta t(H, T_k)[\bar f] = \delta t(H, T_k)[f].
\]
\end{lemma}

\begin{proof}
The cone $\mathcal F_k$ is defined by pointwise inequalities and a linear
constraint, all of which are $G_k$-invariant.  Hence $G_k$ preserves
$\mathcal F_k$, and so does the averaging operation, which is a convex
combination of $G_k$-images.

The second statement uses Lemma~\ref{lem:K-block-constant}: $\mathcal K$ is
$G_k$-invariant.  Therefore $\int \mathcal K\, f\dd x\dd y$ is unchanged when $f$
is replaced by any $G_k$-image, and hence by the average $\bar f$.
\end{proof}

\begin{lemma}[Symmetric perturbations form the $S_{k,a}$ tangent line]
\label{lem:symmetric-perturbations}
Every $G_k$-invariant $\bar f\in\mathcal F_k$ has the form
\[
   \bar f(x,y) = M_w \chi_{\mathrm{within}}(x,y) - M_c \chi_{\mathrm{cross}}(x,y)
\]
for some scalars $M_w\ge 0$ and $M_c\ge 0$, with $\chi_\bullet$ the indicators
of the within-part and cross-part regions.  Moreover the edge-density
constraint $\int \bar f = 0$ forces
\[
   M_w \cdot \mathrm{vol}(\mathrm{within}) = M_c \cdot \mathrm{vol}(\mathrm{cross}),
   \quad\text{i.e.\ }\quad
   \tfrac{M_w}{k} = M_c\cdot \tfrac{k-1}{k},
   \quad
   M_c = \tfrac{M_w}{k-1}.
\]
\end{lemma}

\begin{proof}
A $G_k$-invariant symmetric function on $[0,1]^2$ must, by the same argument as
in Lemma~\ref{lem:K-block-constant}, be constant on within-part cells and
constant on cross-part cells.  The feasibility inequalities then translate into
$\bar f|_{\mathrm{within}}\ge 0$ and $\bar f|_{\mathrm{cross}}\le 0$, i.e.\ the
indicated $M_w, M_c\ge 0$.  Volumes:
$\mathrm{vol}(\mathrm{within}) = k\cdot(1/k)^2 = 1/k$ and
$\mathrm{vol}(\mathrm{cross}) = 1 - 1/k$.  The constraint $\int \bar f = 0$
yields the stated relation.
\end{proof}

\begin{corollary}[Symmetric feasibility cone is one-dimensional]
\label{cor:1d-cone}
The $G_k$-invariant subset of $\mathcal F_k$ is the half-line
\[
   \{M_w\cdot \chi_{\mathrm{within}} - \tfrac{M_w}{k-1}\cdot\chi_{\mathrm{cross}}
   : M_w\ge 0\}.
\]
Up to the positive scalar $M_w$ this is the tangent direction of the family
$\{S_{k,a}\}_{a\ge 0}$ at $a=0$ (note $S_{k,0}=T_k$).
\end{corollary}

\begin{proof}
$\frac{d}{da}S_{k,a}|_{a=0}(x,y)$ equals $+1$ on within-part cells and
$-1/(k-1)$ on cross-part cells.  This is the case $M_w=1$.
\end{proof}

\begin{theorem}[First-order stability of $T_k$ for $k\ge 6$]
\label{thm:first-order-stability}
Let $H$ be the six-vertex Problem-1 graph.  Then for every $k\ge 6$ and every
$f\in\mathcal F_k$,
\[
   \delta t(H,T_k)[f]
   = \frac{6k^3-55k^2+142k-111}{k^5}\cdot
   \int_{\mathrm{within}} f \dd x\dd y
   \ \ge\ 0,
\]
with equality iff the symmetrisation $\bar f$ vanishes a.e.\ (equivalently, iff
$M_w(\bar f) = 0$).  In particular $T_k$ is first-order locally stable against
every edge-density-preserving feasible perturbation.
\end{theorem}

\begin{proof}
By Lemma~\ref{lem:averaging}, $\delta t(H,T_k)[f] = \delta t(H,T_k)[\bar f]$, so
we may replace $f$ by its $G_k$-symmetrisation $\bar f$.  By
Corollary~\ref{cor:1d-cone}, $\bar f = M_w\cdot \frac{d}{da}S_{k,a}|_{a=0}$ for
some $M_w\ge 0$.  Hence
\[
   \delta t(H,T_k)[\bar f]
   = M_w\cdot \frac{d}{da}\,t(H,S_{k,a})\Big|_{a=0}
   = M_w\cdot \frac{6k^3-55k^2+142k-111}{k^5},
\]
the derivative being the one computed in \S\ref{ssec:Ska} and
Appendix~\ref{app:Ska}.  The polynomial $6k^3-55k^2+142k-111$ evaluates to $57$
at $k=6$ and has derivative $18k^2-110k+142$, which is positive at $k=5$ and
strictly increasing for $k\ge 5$.  Hence $6k^3-55k^2+142k-111$ is strictly
increasing for $k\ge 5$ and strictly positive for all $k\ge 6$.

Finally, $M_w\cdot \mathrm{vol}(\mathrm{within}) = M_w/k =
\int_{\mathrm{within}}\bar f$, and the symmetrisation is mass-preserving on the
within region, giving the displayed identity.  Strict inequality holds whenever
$M_w(\bar f)>0$, i.e.\ iff $\bar f\not\equiv 0$ as a graphon.
\end{proof}

\begin{remark}[Why this does not settle Problem~1 for $k\ge 6$]
Theorem~\ref{thm:first-order-stability} is a \emph{negative} result for the
$k\ge 6$ part of Problem~1: it says the natural infinitesimal attack cannot
prove that $T_k$ is not the global minimum, in contrast with $k=3,4,5$ where the
perturbation suffices.  Settling the $k\ge 6$ case therefore requires either
(a) a global construction beating $T_k$ at some $k\ge 6$ (e.g.\ a
higher-order/non-symmetric perturbation, or a different structure such as a
circular construction analogous to the $k=5$ one), or (b) a global optimality
proof for $T_k$ in this range.  A higher-order analysis is the next infinitesimal
step: compute the second derivative of $t(H, S_{k,a,b})$ along a two-parameter
symmetric family (a within-block density $a$ and a single off-diagonal block
density perturbation $b$) at $(a,b)=(0,0)$.  If the resulting Hessian has a
negative eigenvalue in the constrained-feasibility cone, that already settles
non-optimality; otherwise we have a strict local minimum within the
two-parameter class and must broaden the search.
\end{remark}

\begin{remark}[Comparison with $k=3,4,5$]
For $k=3,4,5$, the polynomial $6k^3-55k^2+142k-111$ evaluates to $-18,-39,-26$
respectively, so the derivative is \emph{strictly negative}, giving an explicit
edge-density-preserving feasible direction along which $t(H,W)$ \emph{decreases}
from $T_k$ at first order.  Integrating finitely along $S_{k,a}$ yields the
rational competitors of \S\ref{ssec:Ska}.
\end{remark}

% ============================================================================
\section{Problem 2: the original extension is false}
\label{sec:problem2false}

The second problem asks for a convex/concave classification of $f_F(x)$ for
graphs $F$ with no isolated vertices, no isolated edges, and at least one odd
cycle.  It then asks to extend the classification to all graphs with more than
five vertices.  The extension to all larger graphs is false.

\subsection{A counterexample: two disjoint triangles}

Let $F=K_3\sqcup K_3$.  This graph has six vertices, no isolated vertices, no
isolated edges, and is non-bipartite.  For every graphon $W$,
$t(F,W)=t(K_3,W)^2$, and since $z\mapsto z^2$ is increasing on $[0,\infty)$,
$f_F(x)=f_{K_3}(x)^2$.

On the Tur\'an interval $I_2=[\tfrac12,\tfrac23]$, the triangle-density theorem
gives the usual tripartite minimiser with part sizes $s,\frac{1-s}{2},
\frac{1-s}{2}$, $0\le s\le \frac13$.  The edge and triangle densities are
\[
   x=\frac12+s-\frac32s^2,
   \qquad
   f_{K_3}(x)=\frac32s(1-s)^2,
\]
so $f_{K_3\sqcup K_3}(x)=\left(\frac32s(1-s)^2\right)^2$.  A direct calculation
gives
\[
   \frac{d^2}{dx^2}f_{K_3\sqcup K_3}(x)
   =\frac{9(1-s)^2(4s-1)}{2(3s-1)}.
\]
For $0<s<1/4$, this is positive, while for $1/4<s<1/3$, it is negative.  Hence
$f_{K_3\sqcup K_3}$ is neither convex nor concave on $I_2$.

\subsection{A family of product obstructions}

More generally, let $F_m=\underbrace{K_3\sqcup\cdots\sqcup K_3}_{m}$, $m\ge 2$.
Then $f_{F_m}(x)=f_{K_3}(x)^m$, and using the same parameter $s$ on $I_2$,
$f_{F_m}(x)=\left(\frac32s(1-s)^2\right)^m$.  A calculation gives
\[
   \frac{d^2}{dx^2}f_{F_m}(x)
   =
   \frac{m\left(\frac32s(1-s)^2\right)^m
     \bigl((3m-2)s-(m-1)\bigr)}
        {s^2(1-s)^2(3s-1)}.
\]
The sign changes at $s_0=\frac{m-1}{3m-2}\in\left(0,\frac13\right)$.  Thus all
$F_m$ with $m\ge 2$ have mixed curvature on a single Tur\'an interval.

\begin{status}[Consequence]
Any correct extension of Problem 2 must exclude graphs with multiple independent
non-bipartite pieces.  Merely requiring connectedness is probably not the right
repair, because two odd blocks joined by a bridge should still exhibit a coupled
product mechanism.
\end{status}

% ============================================================================
\section{A repaired structural direction: odd-prime graphs}
\label{sec:oddprime}

The counterexample suggests that the obstruction is not simply disconnectedness
but the presence of multiple odd pieces.

\begin{definition}[Odd block and odd-prime graph]
A block of a graph is a maximal $2$-connected subgraph, with bridges treated as
$K_2$-blocks.  A block is \emph{odd} if it contains an odd cycle.  Let
$b_{\rm odd}(F)$ be the number of odd blocks of $F$.  A graph $F$ is called
\emph{odd-prime} if it is connected, non-bipartite, and $b_{\rm odd}(F)=1$.
\end{definition}

\begin{conjecture}[Odd-prime reduction principle]
Graphs with $b_{\rm odd}(F)\ge 2$ should not be expected to obey a pure
convex/concave dichotomy.  For such graphs, the correct theory should decompose
$f_F$ into products or coupled products of the extremal curves associated with
the odd blocks.  For odd-prime graphs, a convex/concave classification may still
exist, but the extremal templates need not be the same as the clique-density
templates.  The right task is to identify, for each odd-prime family, the
correct Tur\'an-type filling construction and then prove the corresponding lower
bound.
\end{conjecture}

This conjecture is deliberately weaker than the original all-graphs extension.
It says that the clique cases are only one part of the story; non-clique
odd-prime graphs may have different extremal mechanisms.

% ============================================================================
\section{Clique cases: a solved concave subfamily}
\label{sec:clique}

The clique cases are completely understood by the clique-density theorem of
Reiher \cite{ReiherCliqueDensity}.  We record the graphon formulation and the
curvature calculation.

Fix $r\ge 3$.  On $I_k=\left[1-\frac1k,1-\frac1{k+1}\right]$, the
Lov\'asz--Simonovits graphon has $k+1$ parts: one part of size $s$ and $k$ equal
parts of size $(1-s)/k$, where $0\le s\le \frac1{k+1}$.  Its edge density is
\[
   x_k(s)=1-s^2-\frac{(1-s)^2}{k},
\]
and its $K_r$-density is
\[
   g_{r,k}(s)=
   \frac{(k)_{r-1}}{k^{r-1}}(1-s)^{r-1}
   \left(
      \frac{k-r+1}{k}+\frac{(r-1)(k+1)}{k}s
   \right),
   \qquad (k)_{r-1}=k(k-1)\cdots(k-r+2).
\]
If $k+1<r$, then $g_{r,k}=0$.  On the first relevant interval for $K_4$,
$r=4,k=3$, this simplifies to
\begin{equation}
\label{eq:g43}
   g_{4,3}(s)=\tfrac89\,s(1-s)^3,
   \qquad x_3(s)=1-s^2-\tfrac{(1-s)^2}{3},\quad s\in[0,\tfrac14],
\end{equation}
with $g_{4,3}=0$ at $x=2/3$ ($s=0$) and $g_{4,3}=3/32$ at $x=3/4$ ($s=1/4$,
the balanced $4$-partite graphon).  Inverting $x_3$ on $[0,1/4]$ gives
$s=\tfrac14\bigl(1-\sqrt{9-12x}\bigr)$; this closed form is used in
\S\ref{sec:numerics}.

\begin{theorem}[Clique curves are piecewise concave]
\label{thm:clique-concave}
For every $r\ge 3$, $f_{K_r}$ is concave on every Tur\'an interval $I_k$.
\end{theorem}

\begin{proof}
The clique-density theorem identifies $f_{K_r}$ with the Lov\'asz--Simonovits
curve $g_{r,k}(s)$ on $I_k$.  We compute curvature in the parameter $s$.  First,
$x_k'(s)=\frac{2}{k}\bigl(1-(k+1)s\bigr)>0$ in the interior of the interval.
Therefore the sign of $d^2g_{r,k}/dx^2$ is the sign of
$g_{r,k}''(s)x_k'(s)-g_{r,k}'(s)x_k''(s)$.  A direct calculation gives
\[
\begin{split}
&g_{r,k}''(s)x_k'(s)-g_{r,k}'(s)x_k''(s)
 =-
\frac{2r(r-1)(r-2)}{k^2}
\frac{(k)_{r-1}}{k^{r-1}}
(1-s)^{r-3}\bigl(1-(k+1)s\bigr)^2.
\end{split}
\]
This is non-positive.  Hence $f_{K_r}$ is concave on $I_k$.
\end{proof}

\begin{remark}
The clique case is the solved model problem: exact curves, exact extremisers,
and the desired concavity.  The non-clique cases below show that one cannot
simply attach rooted bipartite decorations to clique extremisers and expect the
same templates to remain optimal.
\end{remark}

% ============================================================================
\section{Two non-clique cases: the clique template is not optimal}
\label{sec:nonclique}

Let $P=K_4^\dagger$ denote $K_4$ with one pendant edge attached to one vertex of
the $K_4$, and let $J=K_3\cup_{K_2}K_4$ denote the graph obtained by gluing one
edge of a triangle to one edge of a $K_4$.  Both appear in the original list of
proposed concave cases.  The following calculation shows that the usual
Lov\'asz--Simonovits $4$-partite clique template is not optimal for either graph
at edge density $17/25$.

\subsection{The Lov\'asz--Simonovits value at $x=17/25$}

On $I_3=[2/3,3/4]$, the clique-density template is the complete $4$-partite
graphon with part sizes $s,\frac{1-s}{3},\frac{1-s}{3},\frac{1-s}{3}$.  Its edge
density is $x(s)=\frac23+\frac{2s}{3}-\frac{4s^2}{3}$.  Setting $x=17/25$ gives
$s=\frac{5-\sqrt{21}}{20}$.  For this graphon,
\[
   t(P)=\frac23s(1-s)^3
   =\frac{291}{10000}-\frac{7\sqrt{21}}{2000}
   \approx 0.013060985,
   \qquad
   t(J)=\frac{97}{5000}-\frac{7\sqrt{21}}{3000}
   \approx 0.008707323.
\]

\subsection{A better tripartite filling graphon}

Partition $[0,1]$ into three parts $A\sqcup B\sqcup C$ with
$|A|=\alpha=\frac7{20}$, $|B|=|C|=\frac{13}{40}$.  Put all cross-edges between
distinct parts, put no edges inside $B$ or $C$, and inside $A$ put a $q$-regular
triangle-free graphon with $q=\frac{11}{98}$ (for example, split $A$ into two
equal subparts and put value $2q=11/49$ between the subparts and $0$ inside
each).  The edge density is
\[
   t(K_2,W)=\frac12+\alpha-\frac32\alpha^2+\alpha^2q
   =\frac{533}{800}+\frac{11}{800}=\frac{17}{25}.
\]
For $P$, a direct rooted count gives
\[
   t(P,W)=\frac32\alpha^2(1-\alpha)^2q
          \left(\frac{3-\alpha}{2}+\alpha q\right)
   =\frac{1065207}{89600000}\approx 0.011888471,
\]
strictly smaller than the LS value $0.013060985$.  For $J$, the same graphon
gives
\[
   t(J,W)=\alpha^2(1-\alpha)^2q
          \left(\frac32-\alpha+2\alpha q\right)
       =\frac{79937}{11200000}\approx 0.007137232,
\]
strictly smaller than the LS value $0.008707323$.

\begin{status}[Consequence]
The two non-clique cases $P$ and $J$ cannot be solved by importing the
clique-density extremiser.  Their extremisers, if the natural candidates are
correct, involve a complete tripartite skeleton with a triangle-free filling
inside one part.
\end{status}

% ============================================================================
\section{The tripartite Tur\'an-filling reduction}
\label{sec:filling}

We now derive the exact variational formula inside the natural family suggested
by the previous examples.  This part is rigorous as a reduction to a
finite-dimensional optimisation problem; the remaining global problem (every
minimiser can be chosen from this family) is discussed in \S\ref{sec:global}.

\subsection{The family}

Let $[0,1]=A\sqcup B\sqcup C$ with $|A|=\alpha$, $|B|=|C|=\beta=\frac{1-\alpha}{2}$.
Put all cross-edges between $A,B,C$, no edges inside $B$ or $C$, and a
triangle-free graphon $F$ inside $A$.  Normalise $F$ on a probability space of
measure $1$ and write
\[
   q=t(K_2,F),
   \qquad d_F(u)=\int F(u,v)\dd v,
   \qquad D_2(F)=\int d_F(u)^2\dd u.
\]
By the graphon form of Mantel's theorem, $0\le q\le \frac12$.  The edge density
of the full graphon is $x=\frac12+\alpha-\frac32\alpha^2+\alpha^2q$, so for
fixed $x$,
\[
   q=q_x(\alpha)=\frac{x-\frac12-\alpha+\frac32\alpha^2}{\alpha^2}.
\]
The constraint $0\le q_x(\alpha)\le 1/2$ is, for $x\in[2/3,3/4]$, equivalent to
\[
   \alpha\in A_x:=\left[
      \frac{1-\sqrt{3-4x}}2,
      \frac{1+\sqrt{3-4x}}2
   \right].
\]

\subsection{Counting $P=K_4^\dagger$}

\begin{lemma}[Density of $P$ in the filling family]
For the tripartite filling graphon above,
\[
   t(P,W)=\frac32\alpha^2(1-\alpha)^2
          \left[
             \frac{3-\alpha}{2}q+\alpha D_2(F)
          \right].
\]
Consequently, for fixed $\alpha$ and $q$, the minimum over triangle-free
fillings $F$ is attained by a $q$-regular filling and equals
$\frac32\alpha^2(1-\alpha)^2q\left(\frac{3-\alpha}{2}+\alpha q\right)$.
\end{lemma}

\begin{proof}
Every $K_4$ in the full graphon must use two adjacent vertices inside $A$, one
vertex in $B$, and one vertex in $C$.  The pendant edge may attach to the
distinguished $K_4$ vertex; summing over the four possible positions of the
distinguished vertex gives the displayed formula.  Jensen's inequality gives
$D_2(F)\ge q^2$ with equality exactly for $q$-regular fillings, which exist for
every $q\in[0,1/2]$ (split the space into two equal parts, value $2q$ across).
\end{proof}

Therefore the reduced variational problem for $P$ is
\[
\boxed{
   \Psi_P(x)=
   \min_{\alpha\in A_x}
   \frac32\alpha^2(1-\alpha)^2q_x(\alpha)
   \left(\frac{3-\alpha}{2}+\alpha q_x(\alpha)\right)
   =
   \min_{\alpha\in A_x}
   \frac{3(1-\alpha)^2
     (2\alpha^2+\alpha+2x-1)
     (3\alpha^2-2\alpha+2x-1)}{8\alpha}.
}
\]

\subsection{Counting $J=K_3\cup_{K_2}K_4$}

\begin{lemma}[Density of $J$ in the filling family]
For the tripartite filling graphon above,
\[
   t(J,W)=\alpha^2(1-\alpha)^2
          \left[
             \left(\frac32-\alpha\right)q+2\alpha D_2(F)
          \right],
\]
with in-fill minimum (at a $q$-regular filling)
$\alpha^2(1-\alpha)^2q\left(\frac32-\alpha+2\alpha q\right)$.
\end{lemma}

\begin{proof}
Label the shared edge $1,2$, the remaining triangle vertex $3$, and the two
remaining $K_4$ vertices $4,5$.  The $K_4$ portion uses two adjacent vertices in
$A$, one in $B$, one in $C$.  Summing the rooted contributions over the three
cases (shared edge inside $A$, with one endpoint in $A$, or outside $A$) gives
the formula; the Jensen step is as before.
\end{proof}

Thus the reduced variational problem for $J$ is
\[
\boxed{
   \Psi_J(x)=
   \min_{\alpha\in A_x}
   \alpha^2(1-\alpha)^2q_x(\alpha)
   \left(\frac32-\alpha+2\alpha q_x(\alpha)\right)
   =
   \min_{\alpha\in A_x}
   \frac{(1-\alpha)^2
     (3\alpha^2-2\alpha+2x-1)
     (4\alpha^2-\alpha+4x-2)}{4\alpha}.
}
\]

\subsection{Stationary equations and algebraic branches}
\label{ssec:branches}

For $P$, an interior stationary point satisfies
\[
\begin{split}
0={}&2\alpha^3q^2+9\alpha^3q-3\alpha^3
+2\alpha^2q^2-9\alpha^2q+13\alpha^2
+4\alpha q-13\alpha+3,
\end{split}
\]
with $q=q_x(\alpha)$.  Solving this quadratic for $q$ gives the relevant branch
\[
   q_P(\alpha)=
   \frac{-9\alpha^2+9\alpha-4+
   \sqrt{105\alpha^4-242\alpha^3+153\alpha^2+8\alpha-8}}
   {4\alpha(\alpha+1)},
   \qquad
   q_P(1/3)=0,\ q_P(1/2)=1/2.
\]
It parametrises a candidate reduced optimum by
$x_P(\alpha)=\frac12+\alpha-\frac32\alpha^2+\alpha^2q_P(\alpha)$,
$\alpha\in[1/3,1/2]$, with value
$\Psi_P(x_P(\alpha))=\frac32\alpha^2(1-\alpha)^2q_P(\alpha)
\left(\frac{3-\alpha}{2}+\alpha q_P(\alpha)\right)$.

For $J$, the stationarity equation is
\[
\begin{split}
0={}&4\alpha^3q^2+18\alpha^3q-6\alpha^3
+4\alpha^2q^2-24\alpha^2q+17\alpha^2
+8\alpha q-14\alpha+3,
\end{split}
\]
with relevant branch
\[
   q_J(\alpha)=
   \frac{-(3\alpha-2)^2+
   \sqrt{105\alpha^4-260\alpha^3+204\alpha^2-52\alpha+4}}
   {4\alpha(\alpha+1)},
   \qquad
   q_J(1/3)=0,\ q_J(1/2)=1/2.
\]
At the endpoints,
\[
   \Psi_P(2/3)=0,\quad \Psi_P(3/4)=\frac9{128},
   \qquad
   \Psi_J(2/3)=0,\quad \Psi_J(3/4)=\frac3{64}.
\]

\begin{remark}[Status of the algebraic verification --- now rigorous]
The counting and Jensen reductions above are fully rigorous.  The remaining
one-dimensional minimisation (branch uniqueness, the range
$q_P,q_J\in(0,1/2)$, monotonicity of $x_P,x_J$, the endpoint values, the
curvature, and the strict beating of the feasibility boundary) is certified by
exact Sturm sequences in \S\ref{sec:certificates}, reproduced by
\texttt{scripts/verify\_reduced\_problems.py} and \texttt{verify\_gap.py}.
\end{remark}

\subsection{Numerical example at $x=17/25$}

Solving the reduced problems at $x=\frac{17}{25}$ gives
\[
   \alpha_P\approx 0.3489617543,\quad q_P\approx 0.1125007820,\quad
   \Psi_P(17/25)\approx 0.01188713928,
\]
and
\[
   \alpha_J\approx 0.3509647872,\quad q_J\approx 0.1120315971,\quad
   \Psi_J(17/25)\approx 0.007136528755.
\]
The hand-picked values $\alpha=7/20$ and $q=11/98$ of \S\ref{sec:nonclique} are
therefore very close to the reduced optima but not exactly optimal.

% ============================================================================
\section{Rigorous certificates for $\Psi_P,\Psi_J$ on $[2/3,3/4]$}
\label{sec:certificates}

This section records exact algebraic certificates for the reduced problems
$\Psi_P(x)=\min_\alpha\Phi_P(\alpha,x)$, $\Psi_J(x)=\min_\alpha\Phi_J(\alpha,x)$
of \S\ref{sec:filling}.  Everything is verified by
\texttt{scripts/verify\_reduced\_problems.py} (items in
\S\ref{ssec:cert-branch}--\S\ref{ssec:cert-curvature}) and
\texttt{verify\_gap.py} (\S\ref{ssec:cert-gap}) via \texttt{sympy}.  The
substantive new finding is the mixed curvature
(Theorem~\ref{thm:mixed-curvature}).

\subsection{Setup and key polynomials}

For $\alpha\in(0,1)$ we have the edge-density constraint
\begin{equation}
\label{eq:edgedensity}
   x = \tfrac12 + \alpha - \tfrac32\alpha^2 + \alpha^2 q,
   \qquad
   q_x(\alpha) = \frac{x - \tfrac12 - \alpha + \tfrac32\alpha^2}{\alpha^2},
\end{equation}
and the $(\alpha,q)$ rooted-counting forms
\begin{align}
\Phi_P(\alpha, q) &= \tfrac32\alpha^2(1-\alpha)^2 q\!\left(\tfrac{3-\alpha}{2}+\alpha q\right),
\label{eq:PhiPaq}
\\
\Phi_J(\alpha, q) &= \alpha^2(1-\alpha)^2 q\!\left(\tfrac32-\alpha+2\alpha q\right).
\label{eq:PhiJaq}
\end{align}
After substituting $q = q_x(\alpha)$, the explicit closed forms
\begin{align}
\Phi_P(\alpha, x) &= \frac{3(1-\alpha)^2(2\alpha^2+\alpha+2x-1)(3\alpha^2-2\alpha+2x-1)}{8\alpha},
\label{eq:PhiPx}
\\
\Phi_J(\alpha, x) &= \frac{(1-\alpha)^2(3\alpha^2-2\alpha+2x-1)(4\alpha^2-\alpha+4x-2)}{4\alpha}
\label{eq:PhiJx}
\end{align}
agree with~\eqref{eq:PhiPaq}--\eqref{eq:PhiJaq} via~\eqref{eq:edgedensity}
(symbolically verified).  The stationary equations
$\partial_\alpha \Phi_P = \partial_\alpha\Phi_J = 0$ along the edge-density
constraint reduce (after clearing $(1-\alpha)$) to
\begin{equation}
\label{eq:statPJ}
   A_*(\alpha) q^2 + B_*(\alpha) q + C_*(\alpha) = 0,
\end{equation}
where
\begin{align*}
A_P &= 2\alpha^2(\alpha+1), &
A_J &= 4\alpha^2(\alpha+1),
\\
B_P &= \alpha(9\alpha^2-9\alpha+4), &
B_J &= 2\alpha(3\alpha-2)^2,
\\
C_P &= -3(\alpha-\tfrac13)(\alpha-1)(\alpha-3), &
C_J &= -3(\alpha-\tfrac13)(\alpha-1)(2\alpha-3),
\end{align*}
with discriminants
\begin{align}
\label{eq:DeltaP}
\Delta_P(\alpha) &= 105\alpha^4 - 242\alpha^3 + 153\alpha^2 + 8\alpha - 8,
\\
\label{eq:DeltaJ}
\Delta_J(\alpha) &= 105\alpha^4 - 260\alpha^3 + 204\alpha^2 - 52\alpha + 4,
\end{align}
in the sense that $B_P^2 - 4 A_P C_P = \alpha^2\Delta_P$ and
$B_J^2 - 4 A_J C_J = 4\alpha^2\Delta_J$.  (Note the doubled $\sqrt\Delta_J$
prefactor for $J$.)

\subsection{Discriminant positivity, branch uniqueness, and range}
\label{ssec:cert-branch}

\begin{lemma}[Discriminants are positive]
\label{lem:disc-positive}
$\Delta_P(\alpha) > 0$ and $\Delta_J(\alpha) > 0$ on $(1/3,1/2)$.
\end{lemma}

\begin{proof}
Each $\Delta$ is a quartic with rational coefficients; its Sturm sequence shows
zero real roots in $(1/3,1/2)$.  Endpoint values $\Delta_P(1/3)=4$,
$\Delta_P(1/2)=169/16$, $\Delta_J(1/3)=1$, $\Delta_J(1/2)=49/16$ are positive.
\end{proof}

Hence both stationary quadratics have two real roots
\[
   q^{\pm}_P = \frac{-B_P\pm\alpha\sqrt{\Delta_P}}{2A_P},
   \quad
   q^{\pm}_J = \frac{-B_J\pm 2\alpha\sqrt{\Delta_J}}{2A_J},
\]
and we denote the larger ($+$ sign) root by $q_P$, $q_J$.

\begin{theorem}[Unique stationary branch with $q\in(0,1/2)$]
\label{thm:branch-range}
For every $\alpha\in(1/3,1/2)$, $0 < q_P(\alpha) < \tfrac12$ and
$0 < q_J(\alpha) < \tfrac12$, the two roots of~\eqref{eq:statPJ} have opposite
signs, and $q_P,q_J$ are the unique stationary branches with $q\in(0,1/2)$.
Moreover $q_P(\tfrac13)=q_J(\tfrac13)=0$ and $q_P(\tfrac12)=q_J(\tfrac12)=\tfrac12$.
\end{theorem}

\begin{proof}
On $(1/3,1/2)$, $A_P,A_J>0$ (Sturm).  The constant term
$C_P=-3(\alpha-1/3)(\alpha-3)(\alpha-1)$ has factor signs $+,-,-$, so $C_P<0$;
similarly $C_J=-3(\alpha-1/3)(\alpha-1)(2\alpha-3)<0$.  Since $A>0$, $C<0$, the
product of roots $C/A<0$, so the roots have opposite sign and the ``$+$'' root is
the unique positive one.  For the upper bound, evaluate the stationary quadratic
at $q=1/2$:
\begin{align*}
A_P/4 + B_P/2 + C_P &= 2(\alpha-1/2)(\alpha^2+5\alpha-3), \\
A_J/4 + B_J/2 + C_J &= 2(\alpha-1/2)(2\alpha^2+4\alpha-3).
\end{align*}
On $(1/3,1/2)$, $(\alpha-1/2)<0$; and $\alpha^2+5\alpha-3$ (value $-11/9$ at
$1/3$, $-1/4$ at $1/2$, increasing) and $2\alpha^2+4\alpha-3$ (value $-13/9$ at
$1/3$, $-1/2$ at $1/2$, increasing) are both negative.  Hence each product is
positive at $q=1/2$; with $A>0$ this forces both roots $<1/2$.  Boundary values
are direct evaluations using the endpoint discriminants.
\end{proof}

\subsection{Monotonicity and the parametrisation of $[2/3,3/4]$}
\label{ssec:cert-mono}

\begin{theorem}[Monotonicity]
\label{thm:monotonicity}
$x_P'(\alpha) > 0$ and $x_J'(\alpha) > 0$ on $(1/3,1/2)$.  Combined with
Theorem~\ref{thm:branch-range}, $\alpha\mapsto x_P(\alpha)$ and
$\alpha\mapsto x_J(\alpha)$ are continuous strictly increasing bijections from
$[1/3,1/2]$ onto $[2/3,3/4]$.
\end{theorem}

\begin{proof}
Differentiating $x_P(\alpha)=\tfrac12+\alpha-\tfrac32\alpha^2+\alpha^2 q_P$ and
using~\eqref{eq:statPJ} for implicit differentiation,
\[
   x_P'(\alpha)
   = \frac{(1-3\alpha+2\alpha q)(2A_P q + B_P) - \alpha^2(A_P' q^2 + B_P' q + C_P')}
          {2 A_P q + B_P},
\]
whose denominator equals $\alpha\sqrt{\Delta_P}>0$ on the $+$ branch
(Lemma~\ref{lem:disc-positive}).  Reducing $q^2$ modulo the stationary quadratic,
the numerator is linear in $q$, of the form $L_P(\alpha)+M_P(\alpha)q$ with (after
clearing $\alpha+1$)
\[
L_P = -15\alpha^5 - 15\alpha^4 - 11\alpha^3 + 19\alpha^2 - 2\alpha,
\quad
M_P = -30\alpha^5 - 38\alpha^4 + 14\alpha^3.
\]
On the $+$ branch the numerator becomes
$P_{\mathrm{lin}}(\alpha) + P_{\mathrm{sq}}(\alpha)\sqrt{\Delta_P}$ with
\[
P_{\mathrm{lin}} = 210\alpha^8 - 48\alpha^7 - 452\alpha^6 + 310\alpha^5 + 12\alpha^4 - 8\alpha^3,
\quad
P_{\mathrm{sq}}  = -30\alpha^6 - 38\alpha^5 + 14\alpha^4.
\]
A Sturm check gives $P_{\mathrm{lin}}>0$ and
$P_{\mathrm{lin}}^2 - P_{\mathrm{sq}}^2\Delta_P>0$ on $(1/3,1/2)$ (endpoint
values $1269760/4782969$ at $1/3$, $3087/2048$ at $1/2$).  Hence
$|P_{\mathrm{lin}}|>|P_{\mathrm{sq}}\sqrt{\Delta_P}|$, and with
$P_{\mathrm{lin}}>0$ this gives $P_{\mathrm{lin}}+P_{\mathrm{sq}}\sqrt{\Delta_P}>0$
regardless of the sign of $P_{\mathrm{sq}}$.  Thus $x_P'>0$.  The same argument
applies to $J$ with
\begin{align*}
L_J &= -30\alpha^5+15\alpha^4+2\alpha^3-\alpha^2+2\alpha,
&
M_J &= -60\alpha^5-70\alpha^4+40\alpha^3,
\\
P_{\mathrm{lin}}^{(J)} &= 840\alpha^8 - 300\alpha^7 - 1784\alpha^6 + 1528\alpha^5 - 312\alpha^4 + 16\alpha^3,
&
P_{\mathrm{sq}}^{(J)}  &= -120\alpha^6 - 140\alpha^5 + 80\alpha^4,
\end{align*}
with $P_{\mathrm{lin}}^{(J)}>0$ and
$(P_{\mathrm{lin}}^{(J)})^2-(P_{\mathrm{sq}}^{(J)})^2\Delta_J>0$ (endpoints
$1275904/4782969$ at $1/3$, $99/16$ at $1/2$).  Continuity on $[1/3,1/2]$ and the
endpoint values $x_*(1/3)=2/3$, $x_*(1/2)=3/4$ give the bijections.
\end{proof}

\subsection{Endpoint values}

\begin{theorem}[Endpoint values]
\label{thm:endpoints}
$\Psi_P(2/3) = \Psi_J(2/3) = 0$, $\Psi_P(3/4) = 9/128$, $\Psi_J(3/4) = 3/64$.
\end{theorem}

\begin{proof}
At $\alpha=1/3$, $q_P=q_J=0$, so \eqref{eq:PhiPaq}--\eqref{eq:PhiJaq} vanish.  At
$\alpha=1/2$, $q_P=q_J=1/2$:
$\Phi_P(\tfrac12,\tfrac12)=\tfrac32\cdot\tfrac14\cdot\tfrac14\cdot\tfrac12\cdot
(\tfrac{5/2}{2}+\tfrac14)=9/128$ and
$\Phi_J(\tfrac12,\tfrac12)=\tfrac14\cdot\tfrac14\cdot\tfrac12\cdot
(\tfrac32-\tfrac12+\tfrac12)=3/64$.  By Theorem~\ref{thm:monotonicity} these
endpoints correspond to $x=2/3$ and $x=3/4$, and the parametrisation points are
unique.
\end{proof}

\subsection{Mixed curvature: the new finding}
\label{ssec:cert-curvature}

At an interior critical point $\alpha^*=\alpha^*(x)$ of
$\alpha\mapsto\Phi(\alpha,x)$, the envelope formula
$\Psi''(x) = \det\nabla^2_{(\alpha,x)}\Phi/\Phi_{\alpha\alpha}$ combined with
$\Phi_{\alpha\alpha}>0$ gives
$\operatorname{sign}\Psi''(x)=\operatorname{sign}\det\nabla^2\Phi$.

\begin{proposition}[$\Phi_{\alpha\alpha}>0$ on the branch]
\label{prop:Phiaa-positive}
For each of $\Phi=\Phi_P,\Phi_J$, the restriction
$\Phi_{\alpha\alpha}(\alpha,x_*(\alpha))$ is strictly positive on
$\alpha\in(1/3,1/2)$.
\end{proposition}

\begin{proof}
After eliminating $x$, the restriction is $L(\alpha)+M(\alpha)\sqrt{\Delta_*}$
over a common denominator $\propto\alpha(\alpha+1)^2>0$, with
\begin{align*}
\text{(P)}\quad L_{P,\mathrm{aa}} &= -315\alpha^6+387\alpha^5+606\alpha^4-1143\alpha^3+447\alpha^2+30\alpha-12,\\
            M_{P,\mathrm{aa}} &= 45\alpha^4+12\alpha^3-78\alpha^2+21\alpha,\\
\text{(J)}\quad L_{J,\mathrm{aa}} &= -210\alpha^6+285\alpha^5+371\alpha^4-828\alpha^3+460\alpha^2-82\alpha+4,\\
            M_{J,\mathrm{aa}} &= 30\alpha^4+5\alpha^3-55\alpha^2+20\alpha.
\end{align*}
For each, Sturm shows $L>0$ and $L^2-M^2\Delta>0$ on $(1/3,1/2)$ (endpoints
$L_{P,\mathrm{aa}}^2-M_{P,\mathrm{aa}}^2\Delta_P=1269760/6561$ at $1/3$, $27783/512$
at $1/2$; $L_{J,\mathrm{aa}}^2-M_{J,\mathrm{aa}}^2\Delta_J=318976/59049$ at $1/3$,
$99/16$ at $1/2$).  As in Theorem~\ref{thm:monotonicity}, $L+M\sqrt\Delta>0$.
\end{proof}

The Hessian determinant along the branch is likewise
$L(\alpha)+M(\alpha)\sqrt{\Delta_*}$ over a positive denominator
$\propto\alpha(\alpha+1)^2$, with
\begin{align*}
\text{(P)}\quad L_{P,\det} &= -2835\alpha^7+8280\alpha^6-2385\alpha^5-16236\alpha^4
+21915\alpha^3-9720\alpha^2+729\alpha+252,\\
            M_{P,\det} &= 270\alpha^5-468\alpha^4-342\alpha^3+1134\alpha^2-720\alpha+126,
\\
\text{(J)}\quad L_{J,\det} &= -315\alpha^7+965\alpha^6-403\alpha^5-1753\alpha^4
+2786\alpha^3-1674\alpha^2+432\alpha-38,\\
            M_{J,\det} &= 30\alpha^5-55\alpha^4-35\alpha^3+135\alpha^2-95\alpha+20.
\end{align*}

\begin{theorem}[Mixed curvature]
\label{thm:mixed-curvature}
There exist
\[
   x^*_P \in [0.6787853539,\,0.6787853540],
   \qquad
   x^*_J \in [0.7132424244,\,0.7132424245]
\]
such that $\Psi_P''>0$ on $(2/3,x^*_P)$, $\Psi_P''<0$ on $(x^*_P,3/4)$, and
likewise for $J$ at $x^*_J$.  In particular neither $\Psi_P$ nor $\Psi_J$ is
convex or concave on $[2/3,3/4]$, and each has a unique inflection point inside
that interval.
\end{theorem}

\begin{proof}
By Proposition~\ref{prop:Phiaa-positive},
$\operatorname{sign}\Psi_*''=\operatorname{sign}(L_{*,\det}+M_{*,\det}\sqrt{\Delta_*})$
on the branch.  The number of zeros of $L+M\sqrt\Delta$ on an interval is at most
the number of real roots of $R:=L^2-M^2\Delta$, since every such zero is a zero
of $R=(L+M\sqrt\Delta)(L-M\sqrt\Delta)$.  Sturm gives: $R_J$ has exactly one real
root in $(1/3,1/2)$; $R_P$ has exactly two.

\emph{Existence and bracketing.}  Direct rational evaluation gives
$\det\nabla^2\Phi_P>0$ at $\alpha=1/3+\tfrac1{300}$ and $<0$ at
$\alpha=1/2-\tfrac1{100}$ (and the same for $J$).  Bisecting
$[\tfrac13+10^{-6},\tfrac12-10^{-6}]$ to width $<10^{-12}$ yields explicit
rational brackets with $\alpha^*_P\approx0.347393$, $\alpha^*_J\approx0.401405$;
applying $x_P,x_J$ and Theorem~\ref{thm:monotonicity} gives the stated $x$-brackets
($x^*_P\approx0.67878535394$, $x^*_J\approx0.71324242444$).

\emph{Uniqueness for $J$.}  $R_J$ has exactly one root in $(1/3,1/2)$, so at most
one sign change; existence forces exactly one.

\emph{Uniqueness for $P$.}  $R_P$ has two roots in $(1/3,1/2)$; we must separate
roots of $L_P+M_P\sqrt{\Delta_P}$ from those of the conjugate
$L_P-M_P\sqrt{\Delta_P}$.  If $L+M\sqrt\Delta=0$ then $L,M$ have opposite signs;
if $L-M\sqrt\Delta=0$ they have the same sign.  Sturm isolates the two roots into
$(1/3,2/5)$ and $(2/5,1/2)$, on which:
\begin{center}
\begin{tabular}{l|cc|c}
sub-interval & $\operatorname{sign} L_{P,\det}$ & $\operatorname{sign} M_{P,\det}$ & branch \\ \hline
$(1/3,2/5)$ & $+1$ & $-1$ & $L+M\sqrt\Delta=0$ (true sign change of $\det H$) \\
$(2/5,1/2)$ & $-1$ & $-1$ & $L-M\sqrt\Delta=0$ (NOT a sign change of $\det H$) \\
\end{tabular}
\end{center}
Hence the second Sturm root belongs to the conjugate, and $\det\nabla^2\Phi_P$
has exactly one sign change on $(1/3,1/2)$.
\end{proof}

\begin{remark}[Numerical illustration]
At alpha-values spaced across the branch, $\det\nabla^2\Phi\big|_{\text{branch}}$ is
\[
\begin{array}{c|c}
\alpha & \det H_P \\ \hline
0.3333 & +1.667 \\
0.3400 & +0.846 \\
0.3473 & 0      \\
0.3500 & -0.267 \\
0.4000 & -4.020 \\
0.4900 & -5.711
\end{array}
\quad
\begin{array}{c|c}
\alpha & \det H_J \\ \hline
0.3667 & +1.960 \\
0.3800 & +1.141 \\
0.4014 & 0      \\
0.4300 & -1.177 \\
0.4600 & -1.422 \\
0.4900 & -2.474
\end{array}
\]
exhibiting the single sign change cleanly in both cases.
\end{remark}

\subsection{Stationary minimum vs.\ feasibility boundary}
\label{ssec:cert-gap}

The arguments above identify $(\alpha^*(x),q^*(x))$ as the unique interior
critical point.  To conclude it achieves the in-family minimum $\Psi_*(x)$, we
compare to the boundary $0\le q\le1/2$.  Only $q=1/2$ can be attained for
$x>2/3$; there the values are constant in $x$ once $\alpha$ is fixed:
$\Phi_P(\alpha,\tfrac12)=\tfrac98\alpha^2(1-\alpha)^2$,
$\Phi_J(\alpha,\tfrac12)=\tfrac34\alpha^2(1-\alpha)^2$.  By Vieta on
$\alpha^2-\alpha+(x-1/2)=0$, $\alpha_-\alpha_+=x-1/2$, so the $q=1/2$ boundary
value is $\tfrac98(x-\tfrac12)^2$ for $P$ and $\tfrac34(x-\tfrac12)^2$ for $J$.

\begin{theorem}[Stationary value beats the boundary]
\label{thm:stat-vs-bdry}
For all $x\in[2/3,3/4]$,
\[
   \Psi_P^{\text{stat}}(x) \le \tfrac98 (x-\tfrac12)^2,
   \qquad
   \Psi_J^{\text{stat}}(x) \le \tfrac34 (x-\tfrac12)^2,
\]
with equality only at $x=3/4$.  Hence on $(2/3,3/4)$ the interior critical point
achieves the in-family minimum $\Psi_*(x)$, attained uniquely there.
\end{theorem}

\begin{proof}
Pulling back to the $\alpha$-parametrisation, define the gaps
$G_P=\tfrac98(x-\tfrac12)^2-\Phi_P$, $G_J=\tfrac34(x-\tfrac12)^2-\Phi_J$.  On the
stationary branch, reducing $q^2$ modulo the stationary quadratic gives
$G_*=L_*^{\mathrm{gap}}(\alpha)+M_*^{\mathrm{gap}}(\alpha)\sqrt{\Delta_*}$ over the
positive denominator $128\alpha^4+256\alpha^3+128\alpha^2$, with
\begin{align*}
L_P^{\mathrm{gap}} &= -6\alpha^3(408\alpha^6 - 2227\alpha^5 + 3510\alpha^4 - 2455\alpha^3 + 652\alpha^2 + 24\alpha - 32),
\\
M_P^{\mathrm{gap}} &= 6\alpha^3(40\alpha^4 - 169\alpha^3 + 153\alpha^2 - 52\alpha + 4),
\\
L_J^{\mathrm{gap}} &= -4\alpha^3(816\alpha^6 - 4163\alpha^5 + 7380\alpha^4 - 6260\alpha^3 + 2690\alpha^2 - 534\alpha + 32),
\\
M_J^{\mathrm{gap}} &= 20\alpha^3(16\alpha^4 - 61\alpha^3 + 66\alpha^2 - 28\alpha + 4).
\end{align*}
Sturm gives $L_*^{\mathrm{gap}}>0$ and
$(L_*^{\mathrm{gap}})^2-(M_*^{\mathrm{gap}})^2\Delta_*>0$ on $(1/3,1/2)$ (endpoint
values $123904/59049$ at $1/3$ and $0$ at $1/2$ for $P$; $173056/531441$ at $1/3$
and $0$ at $1/2$ for $J$), so $G_*>0$ on the open interval.  Boundary:
$G_P(\tfrac13,0)=\tfrac1{32}$, $G_J(\tfrac13,0)=\tfrac1{48}$, and
$G_P(\tfrac12,\tfrac12)=G_J(\tfrac12,\tfrac12)=0$.  Uniqueness on $(2/3,3/4)$
follows from $\Phi_{\alpha\alpha}>0$ (Proposition~\ref{prop:Phiaa-positive}) plus
the single critical point (Theorem~\ref{thm:branch-range}) and the strict gap.
\end{proof}

\begin{remark}[Numerical illustration]
At ten evenly spaced $\alpha\in[0.34,0.49]$:
\begin{center}
\begin{tabular}{c||c|c||c|c}
$\alpha$ & $\Psi_P^{\text{stat}}$ & $\tfrac98(x_P-\tfrac12)^2$
         & $\Psi_J^{\text{stat}}$ & $\tfrac34(x_J-\tfrac12)^2$ \\ \hline
0.34 & 0.00519 & 0.03347 & 0.00270 & 0.02214 \\
0.36 & 0.01970 & 0.04008 & 0.01076 & 0.02609 \\
0.40 & 0.04348 & 0.05245 & 0.02571 & 0.03385 \\
0.44 & 0.05973 & 0.06247 & 0.03745 & 0.04060 \\
0.48 & 0.06851 & 0.06883 & 0.04488 & 0.04543 \\
0.49 & 0.06961 & 0.06972 & 0.04602 & 0.04625
\end{tabular}
\end{center}
The gap decreases monotonically to zero as $\alpha\to 1/2$.
\end{remark}

\subsection{Implication for Problem~2}

\begin{corollary}[In-family curvature, conditional on the global reduction]
\label{cor:infamily-curvature}
Assume the global reduction (Conjecture~\ref{conj:global-reduction}): for
$x\in[2/3,3/4]$ the global minimisers of $\min_W t(P,W)$ and $\min_W t(J,W)$
under $t(K_2,W)=x$ can be chosen from the tripartite Tur\'an-filling family.
Then $f_P=\Psi_P$ and $f_J=\Psi_J$ on $[2/3,3/4]$, and by
Theorem~\ref{thm:mixed-curvature} each has mixed convex-then-concave curvature
there, with unique inflection points $x^*_P,x^*_J$.
\end{corollary}

In particular the original Problem~2 statement, listing $P$ and $J$ as piecewise
concave on Tur\'an intervals, fails even \emph{without} the global reduction: it
is wrong already at the level of the reduced one-dimensional curves, on the very
first interval.  This strengthens the qualitative correction of
\S\ref{sec:nonclique} (strict improvement at $x=17/25$) into a quantitative
claim: the curvature itself differs from the conjectured pattern, with a single
sharply localised inflection point in both cases.

% ============================================================================
\section{Toward the global reduction: strategy and structure}
\label{sec:global}

The reductions of \S\ref{sec:filling}--\S\ref{sec:certificates} are
\emph{in-family} statements.  The remaining global problem is:

\begin{conjecture}[Global tripartite-filling reduction]
\label{conj:global-reduction}
For every $x\in[2/3,3/4]$, $f_P(x)=\min\{t(P,W):t(K_2,W)=x\}=\Psi_P(x)$, with the
minimum achieved by a graphon in the tripartite Tur\'an-filling family
\[
   \mathcal T = \biggl\{\ W\ \biggm|\ \exists [0,1] = A\sqcup B\sqcup C,\ |B|=|C|,
   \begin{aligned}
   &W=1\ \text{on cross-block cells},\\
   &W=0\ \text{within }B,\ \text{within }C,\\
   &W|_{A\times A}\ \text{triangle-free}
   \end{aligned}\ \biggr\}.
\]
The analogous statement holds for $J$ and $\Psi_J$.
\end{conjecture}

\subsection{A rooted-degree integral identity}

Both $t(P,\cdot)$ and $t(J,\cdot)$ are integrals of a ``rooted weight'' against a
``local quantity'', which is why their minimisers are not the pure clique-density
minimisers.  For $P=K_4^\dagger$ with the pendant attached to vertex $4$,
\begin{equation}
\label{eq:P-integral}
   t(P,W) = \int_{[0,1]} d_W(x)\,\kappa_4(x)\dd x,
   \qquad
   d_W(x):=\int W(x,y)\dd y,
\end{equation}
where $\kappa_4(x)$ is the rooted $K_4$-density at $x$ (the probability that
three uniformly random partners of $x$ form a $K_3$ with $x$), so
$t(K_4,W)=\int\kappa_4=\E[\kappa_4(X)]$.  Likewise $J$ integrates over the shared
edge:
\begin{equation}
\label{eq:J-integral}
   t(J,W) = \int_{[0,1]^2} W(x,y)\,\tau_3(x,y)\,\tau_4(x,y)\dd x\dd y,
\end{equation}
with $\tau_3,\tau_4$ the rooted $K_3$- and $K_4$-densities at the edge $\{x,y\}$.

\subsection{Why $f_P\ne f_{K_4}$ in general}

The lower bound on $t(P,W)=\int d_W\kappa_4$ at fixed $t(K_2,W)$ is achieved when
$d_W$ and $\kappa_4$ are \emph{anti-correlated}.  In the LS clique template, the
small singleton part has both highest $\kappa_4$ and highest $d_W$ --- positive
correlation, the wrong way for minimising $t(P,W)$.  In the tripartite-filling
family, vertices in $A$ have degree $1-\alpha+\alpha q$ and (because
$W|_{A\times A}$ is triangle-free) suppressed rooted-$K_4$ density, while $B,C$
vertices have degree $(1+\alpha)/2$ and the cross-clique rooted density.  The
construction is the simplest way to lower $\kappa_4$ on the high-degree vertices
and raise it on the low-degree ones.

\subsection{Plan A: stability + symmetrisation}

This route is closest to the clique case (Reiher 2016) and splits into two steps.

\begin{problem}[Step A1: stability of a $3$-blowup skeleton]
\label{prob:A1}
Let $W$ minimise $t(P,W)$ (resp.\ $t(J,W)$) at $t(K_2,W)=x\in[2/3,3/4]$.  Show
$W$ has a ``$3$-blowup skeleton'': $[0,1]=A\sqcup B\sqcup C$ with $W=1$ a.e.\ on
$A\times B, A\times C, B\times C$ and $W=0$ a.e.\ on $B\times B, C\times C$ (the
internal structure inside $A$ unconstrained).
\end{problem}

\begin{problem}[Step A2: symmetrisation to balanced parts]
\label{prob:A2}
With the skeleton fixed, show one may assume $|B|=|C|$ without increasing
$t(P,W)$ (resp.\ $t(J,W)$), and that $W|_{A\times A}$ may be replaced by a
$q$-regular triangle-free filling of the same edge density.
\end{problem}

If A1 and A2 hold, then $f_P\le\Psi_P^{\text{filling-only}}$; combined with the
in-family lower bound $f_P\ge\Psi_P^{\text{filling-only}}$ of
\S\ref{sec:filling}, this gives $f_P=\Psi_P$ on $[2/3,3/4]$.

\paragraph{Step A2 is now closed (corrected form).}  \S\ref{sec:balancing}
proves Step~A2 in the form actually needed --- a \emph{global} (not per-$\alpha$)
minimisation over unbalanced parts.  The naive per-$\alpha$ balancing statement
(at fixed $|A|=\alpha$, balancing $|B|=|C|$ never increases $t(P,W)$) turns out to
be \emph{false} (Proposition~\ref{prop:counter}); the correct global statement
holds (Theorem~\ref{thm:balancing-main}).

\paragraph{Step A1 is the bottleneck.}  It is a clique-density supersaturation
stability statement:

\begin{lemma}[$K_4$-supersaturation stability, conjectural form]
\label{lem:K4-stability}
There is $C(x)>0$ such that for every $W$ with $t(K_2,W)=x\in(2/3,3/4)$,
\[
   \delta_\square\bigl(W,\ \mathcal{T}_{3}^{1}(x)\bigr)
   \le C(x)\cdot \bigl(t(K_4,W) - g_{4,3}(s_x)\bigr)^{1/2},
\]
where $\delta_\square$ is the cut metric, $\mathcal T_3^1(x)$ is the set of
$3$-blowups of edge density $x$, and $g_{4,3}(s_x)$ is the LS $K_4$-density
\eqref{eq:g43} at edge density $x$ (evaluated at $s_x=\tfrac14(1-\sqrt{9-12x})$).
\end{lemma}

\begin{remark}[Status of Lemma~\ref{lem:K4-stability}]
A cut-metric stability statement of this strength for the LS $4$-partite template
at $K_4$-density is not directly in the literature.  Pikhurko--Razborov-type
results \cite{PikhurkoRazborov} cover the equality case at the clique-density
level; the cut-metric version would require adapting the
regularity/decomposition argument used for triangle supersaturation to the
$4$-partite case, or invoking a more recent clique-density stability result.  It
is the most technical ingredient.  \S\ref{sec:numerics} gives numerical evidence
that the bound's mechanism holds.
\end{remark}

\subsection{Plan B: flag-algebra certificates}

A direct flag-algebra SDP produces certified lower bounds $t(P,W)\ge L(x)$ at
finitely many rational $x$.  For Problem~1 the analogous program for $C_5$ at
$p=1-1/k$ was carried out by Bennett--Dudek--Lidick\'y--Pikhurko (2020)
\cite{BennettDudekLidickyPikhurkoC5}; for our $H$ at $p=4/5$ the target
$24349/187500$ is a small-denominator rational, so the SDP is likely tractable.

\begin{problem}[Step B]
\label{prob:B}
Run a degree-$7$ (over $K_7$-flags) SDP to certify $t(H,W)\ge 24349/187500$ for
$W$ with $t(K_2,W)=4/5$.
\end{problem}

For Problem~2 one needs certificates parametric in $x$.  The mixed curvature
(Theorem~\ref{thm:mixed-curvature}) breaks the simplest interpolation: a single
convex certificate cannot work everywhere, so the interpolation must know about
the inflection $x^*$.

\subsection{Plan C: rooted clique-density inequalities}

A more conceptual route develops sharp rooted clique-density inequalities from
\eqref{eq:P-integral}: at fixed marginals $\mu_d,\mu_\kappa$, the integral
$\int d_W\kappa_4$ is minimised by the negative-rearrangement coupling.

\begin{problem}[Step C, qualitative]
\label{prob:C}
Identify the boundary curve of the achievable region in the
$(\E[d_W^2],\E[\kappa_4])$ plane at fixed $t(K_2,W)=x$ on $[2/3,3/4]$, and express
$\Psi_P(x)$ as a rearrangement-style functional restricted to this boundary.
\end{problem}

This plan is more speculative but would also explain the mixed-curvature
phenomenon: the boundary curve has a kink, and $\Psi_P''$ changes sign there.

\subsection{What the mixed-curvature result tells us about the global problem}

The rigorous mixed curvature (Theorem~\ref{thm:mixed-curvature}) constrains what
kind of global lower bound is achievable: a single convex/concave certificate
cannot prove $t(P,W)\ge\Psi_P(x)$ uniformly; a flag-algebra SDP at a single $x$
interpolated linearly is valid only in a neighbourhood bounded by the distance to
$x^*_P$; and a symmetrisation proof must implicitly distinguish the two curvature
regimes (most naturally via the $x$-dependence of $\alpha^*(x)$, whose
bifurcation between ``$A$ small'' and ``$A\to1/2$'' matches the inflection).  This
suggests Plan A is the most natural, as it respects the parametric structure of
the minimiser.

\subsection{Does the global reduction extend beyond $[2/3,3/4]$?}

For $x\in[3/4,4/5]$ the natural analogue is a $4$-blowup with a triangle-free
filling inside one part; for $[4/5,5/6]$ a $5$-blowup with $K_4$-free filling;
etc.  We have not analysed these.  The mixed-curvature result on $[2/3,3/4]$
leaves the structure on $[3/4,4/5]$ as the next open question.

% ============================================================================
\section{Closing Step A2: the in-family reduction to balanced parts}
\label{sec:balancing}

We complete the in-family side of the reduction by dropping the constraint
$|B|=|C|$ and showing the unbalanced freedom never beats the balanced reduced
curve.  Every algebraic step is reproduced by
\texttt{scripts/verify\_balancing\_reduction.py}.

\subsection{Setup: the unbalanced family and its densities}

Let $[0,1]=A\sqcup B\sqcup C$ with $|A|=\alpha,|B|=\beta,|C|=\gamma$,
$\alpha+\beta+\gamma=1$; cross-edges between $A,B,C$; no edges inside $B$ or $C$;
a triangle-free filling $F$ inside $A$ with $q=t(K_2,F)$, $D_2(F)=\int d_F^2$.
Write $s=\beta+\gamma=1-\alpha$, $p=\beta\gamma$.

\begin{lemma}[Densities, unbalanced]
\label{lem:densities}
For the family above,
\begin{align}
   t(K_2,W) &= 2(\alpha\beta+\alpha\gamma+\beta\gamma)+\alpha^2 q
            = 2\alpha s+2p+\alpha^2 q,\\
   t(P,W)   &= 3\alpha^2\beta\gamma\,\bigl[(2\alpha+3s)\,q+2\alpha D_2(F)\bigr],
            \label{eq:tP}\\
   t(J,W)   &= 2\alpha^2\beta\gamma\,\bigl[(\alpha+3s)\,q+4\alpha D_2(F)\bigr].
            \label{eq:tJ}
\end{align}
Both depend on $F$ only through $q$ and $D_2(F)$, and setting $\beta=\gamma=s/2$
recovers the \S\ref{sec:filling} formulas.
\end{lemma}

\begin{proof}
Every $K_4$ uses two $F$-adjacent vertices in $A$, one in $B$, one in $C$ ($B,C$
independent, so $\le1$ vertex each; three vertices in $A$ would be an
$F$-triangle, excluded).  Summing rooted contributions over the position of the
distinguished/shared vertices, keeping $\beta\ne\gamma$, gives
\eqref{eq:tP}--\eqref{eq:tJ}; the Jensen step $D_2\ge q^2$ is unchanged.  The
computation is reproduced symbolically (as a finite stochastic-block-model sum)
in \texttt{verify\_balancing\_reduction.py}, Part~0.
\end{proof}

By Jensen and monotonicity in $D_2$, at fixed $(\alpha,\beta,\gamma,q)$ --- hence
at fixed edge density --- the minimum over fillings is at $D_2=q^2$.  Henceforth
take $D_2=q^2$ and study
\begin{equation}
\label{eq:Treg}
   t(P,W)=3\alpha^2 p\,q\,(2\alpha q+2\alpha+3s),\qquad
   t(J,W)=2\alpha^2 p\,q\,(4\alpha q+\alpha+3s).
\end{equation}

\subsection{The per-$\alpha$ balancing statement is false}

Fix $\alpha,x$; then $s=1-\alpha$ is fixed and $2p+\alpha^2 q$ is fixed, so
increasing $p$ toward its balanced maximum $s^2/4$ forces $q$ down.

\begin{proposition}[Counterexample]
\label{prop:counter}
At $\alpha=\tfrac{11}{20}$, $x=\tfrac23$, balancing \emph{increases} $t(P,W)$:
the balanced configuration ($p=\tfrac{81}{1600}$, $q=\tfrac{169}{726}$) has
$t(P,W)=\tfrac{4074759}{140800000}\approx0.02894$, whereas the unbalanced
configuration with $q=\tfrac12$ ($p=\tfrac{49}{4800}$) has
$t(P,W)=\tfrac{17787}{1280000}\approx0.01390<0.02894$.
\end{proposition}

Thus the natural per-$\alpha$ balancing claim (the tempting form of Step~A2) is
false.  The violation occurs because, for $\alpha>1/2$, the balanced filling
density needed to reach $x$ is small, and shifting weight from the parts into the
filling (toward the Mantel bound $q=1/2$) lowers $t(P,W)$.  Crucially
$\alpha=\tfrac{11}{20}>\tfrac12$ lies outside the range $[\tfrac13,\tfrac12]$
where the balanced optimum sits; such configurations never realise the
\emph{global} in-family minimum, as we now show.

\subsection{The reduction theorem}

Substituting $p=\tfrac12(x-2\alpha s-\alpha^2 q)$ into \eqref{eq:Treg}, the
in-family density at fixed $x$ is
\begin{equation}
\label{eq:T2}
   T_P(\alpha,q)=\tfrac32\alpha^2\,(x-2\alpha s-\alpha^2 q)\,q\,(2\alpha q+3-\alpha),
   \quad
   T_J(\alpha,q)=\alpha^2\,(x-2\alpha s-\alpha^2 q)\,q\,(4\alpha q+3-2\alpha),
\end{equation}
on $R_x=\{(\alpha,q):\alpha\in(0,1),\ 0\le p\le s^2/4,\ 0\le q\le\tfrac12\}$.

\begin{theorem}[In-family reduction to balanced parts]
\label{thm:balancing-main}
For every $x\in(2/3,3/4)$,
$\min_{R_x}T_P=\Psi_P(x)$ and $\min_{R_x}T_J=\Psi_J(x)$, the balanced reduced
values of \S\ref{sec:filling}, attained on the balanced arc $q=q_x(\alpha)$
($\beta=\gamma$, regular filling).
\end{theorem}

\subsubsection*{The feasible region is the balanced--Mantel lens}

In $(\alpha,q)$ coordinates the constraints read
$q\le\bar q(\alpha)=\frac{x-2\alpha s}{\alpha^2}$,
$q\ge\underline q(\alpha)=\frac{x-2\alpha s-s^2/2}{\alpha^2}$, $q\le\tfrac12$.

\begin{lemma}[Tur\'an thresholds]
\label{lem:turan}
For $x>2/3$ and all $\alpha$, $\bar q(\alpha)>\tfrac12\Leftrightarrow
x>2\alpha-\tfrac32\alpha^2$ and $\underline q(\alpha)>0\Leftrightarrow
x>\tfrac12+\alpha-\tfrac32\alpha^2$, and both right-hand inequalities hold since
\[
   \tfrac23-\bigl(2\alpha-\tfrac32\alpha^2\bigr)=\tfrac16(3\alpha-2)^2\ge0,
   \qquad
   \tfrac23-\bigl(\tfrac12+\alpha-\tfrac32\alpha^2\bigr)=\tfrac16(3\alpha-1)^2\ge0.
\]
\end{lemma}

Hence for $x\in(2/3,3/4)$ the constraints $p\ge0$, $q\ge0$ are slack; the active
boundaries are the \emph{balanced arc} $q=\underline q(\alpha)$ and the
\emph{Mantel arc} $q=\tfrac12$, and $R_x$ is the closed lens between them over
$\alpha\in A_x$.  The two perfect squares are precisely why the first Tur\'an
interval begins at $x=2/3$.

\subsubsection*{No interior local minimum: the saddle certificate}

\begin{lemma}[Interior critical points are saddles]
\label{lem:saddle}
For each $x\in(2/3,3/4)$, $T_P$ (resp.\ $T_J$) has $T_{qq}<0$ at every interior
critical point of $R_x$.  Hence no interior critical point is a local minimum,
and $\min_{R_x}T_P$ is attained on $\partial R_x$.
\end{lemma}

\begin{proof}
For fixed $\alpha$, $T_P(\alpha,\cdot)$ is a cubic in $q$ with leading
coefficient $-3\alpha^5<0$; $T_{qq}$ is affine in $q$ and vanishes on the
rational inflection locus
\[
   q_{\mathrm{infl}}^{P}=\frac{5\alpha^2-7\alpha+2x}{6\alpha^2},
   \qquad
   q_{\mathrm{infl}}^{J}=\frac{10\alpha^2-11\alpha+4x}{12\alpha^2},
\]
with $T_{qq}<0$ iff $q>q_{\mathrm{infl}}$.  A critical point with $T_{qq}=0$ would
lie on the inflection locus; substituting $q=q_{\mathrm{infl}}$ into
$\partial_\alpha T,\partial_q T$ gives polynomials $S,U$ whose common zero forces
$\operatorname{Res}_\alpha(S,U)(x)=0$.  An exact computation
(\texttt{verify\_balancing\_reduction.py}, Part~3) shows this degree-$15$
resultant has \emph{no root in $(2/3,3/4)$} for both $P,J$.  So $T_{qq}\ne0$ at
every critical point, its sign is constant along each branch, and at
$x=\tfrac7{10}$ the unique interior critical point has $T_{qq}<0$.  An indefinite
Hessian is not a local minimum; by compactness the minimum is on $\partial R_x$.
\end{proof}

\subsubsection*{The boundary, and Mantel domination}

On the balanced arc, $T(\alpha,q_x(\alpha))=\Phi_P(\alpha,x)$ of
\eqref{eq:PhiPx}, with minimum $\Psi_P(x)$ over $A_x$
(\S\ref{sec:filling}--\S\ref{sec:certificates}).  On the Mantel arc $q=\tfrac12$,
\begin{equation}
\label{eq:M}
   M_P(\alpha):=T_P(\alpha,\tfrac12)=\tfrac94\alpha^2\bigl(x-2\alpha+\tfrac32\alpha^2\bigr),
   \quad
   M_J(\alpha):=T_J(\alpha,\tfrac12)=\tfrac32\alpha^2\bigl(x-2\alpha+\tfrac32\alpha^2\bigr).
\end{equation}
At $q=\tfrac12$ the filling is complete bipartite, so $W$ is the complete
$4$-partite (Lov\'asz--Simonovits) graphon, and \eqref{eq:M} is its $P$- (resp.\
$J$-) density.

\begin{lemma}[Mantel domination]
\label{lem:mantel}
For $x\in(2/3,3/4)$, $\min_{\alpha\in A_x}M_P(\alpha)\ge\Psi_P(x)$, and likewise
for $J$.
\end{lemma}

\begin{proof}[Proof (status)]
$M_P$ is increasing--decreasing--increasing on $A_x$ with interior local minimum
at $\alpha_M^{+}=\tfrac{3+\sqrt{9-12x}}6$; the left endpoint $\alpha_-$ of $A_x$
is a lens corner where $M_P(\alpha_-)=\Phi_P(\alpha_-)\ge\Psi_P(x)$.  Hence it
suffices that $M_P(\alpha_M^{+})\ge\Psi_P(x)$.  Along the Mantel-minimum locus,
$x=3\mu(1-\mu)$ and $M_P(\alpha_M^{+})=\tfrac94\mu^3-\tfrac{27}8\mu^4$ with
$\mu=\alpha_M^{+}$.  This one-dimensional comparison is verified \emph{exactly}
at $98$ rational points dense in $(2/3,3/4)$
(\texttt{verify\_balancing\_reduction.py}, Part~5), with strict margin
$\ge9\cdot10^{-5}$; see \S\ref{sec:balancing-status} for the parametric Sturm
certificate that upgrades this to all $x$.  Note the cheap pointwise bound
$M_P\ge\Phi_P$ at $\alpha_M^{+}$ \emph{fails}, so the comparison genuinely uses
the balanced \emph{minimum}.
\end{proof}

\begin{proof}[Proof of Theorem~\ref{thm:balancing-main}]
By Lemma~\ref{lem:saddle} the minimum is on $\partial R_x=$ (balanced arc)
$\cup$ (Mantel arc).  The balanced arc has minimum $\Psi_P(x)$; the Mantel arc
has minimum $\ge\Psi_P(x)$ by Lemma~\ref{lem:mantel}.  Hence
$\min_{R_x}T_P=\Psi_P(x)$, attained on the balanced arc; same for $J$.
\end{proof}

\subsection{What this does and does not settle}
\label{sec:balancing-status}

Theorem~\ref{thm:balancing-main} closes \textbf{Step~A2} in its corrected form:
within the tripartite-filling family, the unbalanced freedom never beats the
balanced reduced curve, so the in-family minimum is exactly $\Psi_P,\Psi_J$.
Combined with the mixed-curvature result, the in-family picture on the first
Tur\'an interval is complete and rigorous, modulo the one Sturm formalisation
below.  What remains for $f_P=\Psi_P$ globally is \textbf{Step~A1}.

\paragraph{Rigour status.}  Lemma~\ref{lem:densities}, Proposition~\ref{prop:counter},
Lemma~\ref{lem:turan}, and Lemma~\ref{lem:saddle} are exact and uniform in
$x\in(2/3,3/4)$.  Lemma~\ref{lem:mantel} is proved at a dense rational grid; the
remaining task is the parametric certificate that
$M_P(\alpha_M^{+}(x))-\Psi_P(x)>0$ for all $x$, obtained by eliminating $\mu$
(from $3\mu^2-3\mu+x=0$) and the balanced stationary parameter (from
$\partial_\alpha\Phi_P=0$) and checking the resulting univariate polynomial has
no root in $(2/3,3/4)$ --- the same Sturm endgame as the other certificates.  The
resultant factors as $x^6(2x-1)^4(3x-2)^2(4x-3)^2\cdot D_{18}(x)$, whose only root
of $D_{18}$ in $(2/3,3/4)$, at $x\approx0.74948$, is spurious (the true difference
is $\approx4.2\cdot10^{-5}>0$ there); finishing requires isolating the correct
algebraic branches so the certificate sees only the true difference.

\paragraph{A conceptual reading.}  The Mantel arc is the complete $4$-partite LS
template; Lemma~\ref{lem:mantel} says \emph{the filling construction strictly
beats the clique template} on $(2/3,3/4)$ --- exactly the phenomenon
\S\ref{sec:nonclique} discovered numerically.  Step~A2 thus has clean geometric
content: the only two competitive ways to fill a $3$-blowup to edge density $x$
are (i) balance the parts and put a regular triangle-free filling inside $A$ (the
$\Psi$-construction), or (ii) saturate the filling to a $4$-partite graphon; and
(i) always wins.

% ============================================================================
\section{Numerical evidence for Step A1: a direct step-graphon optimiser}
\label{sec:numerics}

Step~A1 (Problem~\ref{prob:A1}) --- that an \emph{arbitrary} minimiser collapses
onto a $3$-blowup skeleton --- is the remaining bottleneck.  As direct evidence,
\texttt{scripts/optimize\_graphon.py} minimises $t(P,W)$ and (separately)
$t(K_4,W)$ over symmetric step graphons $W$ at fixed edge density $x$, by
projected (Adam) gradient with random restarts.  Densities and gradients are
exact (\texttt{einsum}); the gradient is finite-difference-checked to
$\approx10^{-11}$.

Two quantities are tracked at the achieved density $x'=t(K_2,W_{\mathrm{opt}})$
(to remove the penalty-constraint confound):
\begin{itemize}[itemsep=2pt]
\item \texttt{conj-ok}: whether $t(P,W_{\mathrm{opt}})\ge\Psi_P(x')$, i.e.\
whether the global-min conjecture (Conjecture~\ref{conj:global-reduction}) holds
at the found minimiser;
\item the $K_4$-gap $t(K_4,W_{\mathrm{opt}})-g_{4,3}(x')$, where $g_{4,3}$ is the
LS $K_4$-curve \eqref{eq:g43}.  A small gap is exactly the supersaturation
mechanism of Lemma~\ref{lem:K4-stability}: a near-minimal $K_4$-density forces
$W$ close to a $3$-blowup.
\end{itemize}

\begin{center}
\begin{tabular}{c|c|c|c|c|c|c|c}
$x$ & $x'$ & $\min t(P)$ & $\Psi_P(x')$ & \texttt{conj-ok} & $t(K_4)$@$P$-min & $g_{4,3}(x')$ & $K_4$-gap \\ \hline
0.680 & 0.67864 & 0.011563 & 0.010669 & true & 0.017096 & 0.015667 & 0.00143 \\
0.700 & 0.69882 & 0.033076 & 0.028612 & true & 0.047229 & 0.040680 & 0.00655 \\
0.720 & 0.71872 & 0.050275 & 0.045854 & true & 0.069638 & 0.063409 & 0.00623 \\
\end{tabular}
\end{center}

Two findings:
\begin{enumerate}[label=(\arabic*),itemsep=3pt]
\item \texttt{conj-ok} holds at all three densities: the numerically found
$t(P)$-minimiser obeys $t(P,W_{\mathrm{opt}})\ge\Psi_P(x')$, consistent with
$f_P=\Psi_P$.  (It exceeds $\Psi_P$ by a few $\times10^{-3}$, as expected for a
finite step graphon with a penalised constraint.)
\item The $K_4$-gap is small and positive ($\sim10^{-3}$): the $t(P)$-minimiser
sits just above the $K_4$-minimum, the regime in which the Step~A1 stability
mechanism would force proximity to a $3$-blowup.  Inspecting the sorted degree
profile of $W_{\mathrm{opt}}$ shows a clean two-level structure --- a low-degree
cluster (the $A$-part) and a high-degree cluster pinned near the cross-block
value (e.g.\ six blocks at degree $0.75$ at $x=0.72$) --- i.e.\ exactly the
$3$-blowup skeleton, with the nontrivial structure confined to the low-degree
$A$-part.
\end{enumerate}

This is evidence, not proof: it does not establish the cut-metric stability of
Lemma~\ref{lem:K4-stability}, but it confirms that the conjectured minimiser
structure and the supersaturation mechanism behave as Plan~A predicts on the
first Tur\'an interval.

% ============================================================================
\section{Where the project currently stands}
\label{sec:status}

\subsection{Solved or essentially settled}

\begin{enumerate}[label=(\arabic*),itemsep=2pt]
\item Problem 1: balanced $k$-partite and constant values are explicit.
\item The $k=5$ circular graphon has density $24349/187500$ and beats both
benchmarks; explicit rational competitors beat both benchmarks for $k=3,4,5$.
\item For $k\ge 6$, $T_k$ is first-order locally stable
(Theorem~\ref{thm:first-order-stability}, with the averaging step explicit), so
any all-$k$ non-optimality proof must be global.
\item The original all-graphs extension of Problem 2 is false; $K_3\sqcup K_3$
has mixed curvature on $[1/2,2/3]$.
\item All clique curves $f_{K_r}$ are piecewise concave
(Theorem~\ref{thm:clique-concave}).
\item $K_4^\dagger$ and $K_3\cup_{K_2}K_4$ are \emph{not} minimised by the LS
$4$-partite clique template at $x=17/25$.
\item Inside the tripartite Tur\'an-filling family, $P$ and $J$ reduce to the
explicit one-dimensional curves $\Psi_P,\Psi_J$, which are now certified rigorously:
unique stationary branch (Theorem~\ref{thm:branch-range}), monotone
parametrisation (Theorem~\ref{thm:monotonicity}), endpoints
(Theorem~\ref{thm:endpoints}), the in-family minimum beats the boundary
(Theorem~\ref{thm:stat-vs-bdry}), and \textbf{mixed convex-then-concave
curvature} (Theorem~\ref{thm:mixed-curvature}).
\item \textbf{Step~A2} (in-family reduction to balanced parts) is closed in its
corrected global form (Theorem~\ref{thm:balancing-main}), modulo one Sturm
endgame (\S\ref{sec:balancing-status}).  The naive per-$\alpha$ version is false
(Proposition~\ref{prop:counter}).
\end{enumerate}

\subsection{Main open questions}

\begin{problem}[Problem 1, $k=5$ lower bound]
Prove or disprove $t(H,W)\ge 24349/187500$ whenever $t(K_2,W)=4/5$, and identify
all minimisers, or find a better graphon.
\end{problem}

\begin{problem}[Problem 1, all $k$]
For every $k>2$ decide whether the minimiser at $1-1/k$ is neither $T_k$ nor
constant.  For $k\ge6$ this needs a global construction or a global optimality
proof.
\end{problem}

\begin{problem}[Odd-block decomposition]
Develop a decomposition theory for $f_F$ when $F$ has multiple odd blocks; $mK_3$
shows products of odd pieces destroy any simple dichotomy.
\end{problem}

\begin{problem}[Global reduction for $P$ and $J$ --- Step~A1]
Prove the cut-metric stability of Lemma~\ref{lem:K4-stability} (or otherwise
establish the $3$-blowup skeleton), which with Step~A2
(Theorem~\ref{thm:balancing-main}) yields $f_P=\Psi_P$, $f_J=\Psi_J$ on
$[2/3,3/4]$.  \S\ref{sec:numerics} gives supporting evidence.
\end{problem}

\begin{problem}[Curvature on later Tur\'an intervals]
Determine the analogue of $\Psi_P,\Psi_J$ on $[3/4,4/5]$ ($4$-blowup with a
triangle-free filling) and its curvature signature.
\end{problem}

% ============================================================================
\section{Research directions}
\label{sec:directions}

\paragraph{Flag-algebra certificates.}  The most direct path to the missing lower
bounds.  For Problem 1, certify $t(H,W)-24349/187500\ge0$ under $t(K_2,W)=4/5$
(Problem~\ref{prob:B}).  For Problem 2, certify $t(P,W)\ge\Psi_P(x)$ and
$t(J,W)\ge\Psi_J(x)$ on $[2/3,3/4]$; these are parameter-dependent, and the mixed
curvature forces the interpolation to split at $x^*$.

\paragraph{Symmetrisation and stability (Plan A).}  Prove the $3$-blowup skeleton
(Step~A1, Lemma~\ref{lem:K4-stability}) by adapting $K_4$-supersaturation
stability; Step~A2 is done.  This respects the parametric structure of the
minimiser and is the most natural route.

\paragraph{Rooted density inequalities (Plan C).}  Develop sharp rooted
clique-density inequalities under an edge-density constraint, using
\eqref{eq:P-integral}; this would also explain the mixed-curvature kink
(Problem~\ref{prob:C}).

\paragraph{Formal algebraic verification.}  The one-dimensional certificates of
\S\ref{sec:certificates} are exact and already machine-checked in \texttt{sympy};
a Lean/Sage port would make the finite-dimensional part fully publication-ready.

\paragraph{Revisiting the small-graph classification.}  Because $K_4^\dagger$ and
$K_3\cup_{K_2}K_4$ do not use the clique template, re-enumerate non-bipartite
graphs on $\le5$ vertices, generate candidate Tur\'an-filling templates, solve the
reduced problems exactly, and certify by flag algebras.

% ============================================================================
\section{A compact checklist for collaborators}
\label{sec:checklist}

\begin{enumerate}[label=\textbf{Task \arabic*.},itemsep=3pt]
\item Certify the $k=5$ circular value $24349/187500$ at $4/5$ by flag algebras
(Problem~\ref{prob:B}).
\item Search for global competitors to $T_k$ for $k\ge6$ (local perturbations
will not find them; Theorem~\ref{thm:first-order-stability}); the next
infinitesimal step is the second-order analysis of $t(H,S_{k,a,b})$.
\item Treat $K_3\sqcup K_3$ as the canonical obstruction to the original
Problem 2 extension.
\item Do not assume clique-decorated graphs are minimised by clique-density
templates --- $K_4^\dagger$ and $K_3\cup_{K_2}K_4$ disprove that.
\item Finish the all-$x$ Mantel-domination Sturm certificate
(\S\ref{sec:balancing-status}) to fully close Step~A2.
\item Prove Step~A1: the cut-metric stability of Lemma~\ref{lem:K4-stability}
(literature survey of Pikhurko--Razborov-type results is the first move).
\item Optionally, sample a flag-algebra SDP for $\Psi_P$ at $x=17/25$ to gauge
tractability.
\end{enumerate}

% ============================================================================
\appendix

\section{Exact finite-sum checker for the family $S_{k,a}$}
\label{app:Ska}

\[
   t(H,S_{k,a})=\frac1{k^6}\sum_{c_1,\ldots,c_6\in[k]}
        \prod_{ij\in E(H)}
        \left(a\1\{c_i=c_j\}+\left(1-\frac{a}{k-1}\right)\1\{c_i\ne c_j\}\right).
\]
This finite sum is the source of the rational values in \S\ref{ssec:Ska} and of
\[
   \left.\frac{d}{da}t(H,S_{k,a})\right|_{a=0}
   =\frac{6k^3-55k^2+142k-111}{k^5}.
\]

\section{Derivative formulae for the reduced $P$ and $J$ problems}
\label{app:derivatives}

For $P$, with
$\Phi_P(\alpha,x)=\frac{3(1-\alpha)^2(2\alpha^2+\alpha+2x-1)(3\alpha^2-2\alpha+2x-1)}{8\alpha}$,
\[
\frac{\partial \Phi_P}{\partial \alpha}
=\frac{3(\alpha-1)}{8\alpha^2}N_P(\alpha,x),
\quad
\begin{aligned}
N_P={}&30\alpha^5-22\alpha^4+30\alpha^3x-19\alpha^3
-14\alpha^2x+9\alpha^2\\
&+4\alpha x^2-4\alpha x+\alpha+4x^2-4x+1.
\end{aligned}
\]
For $J$, with
$\Phi_J(\alpha,x)=\frac{(1-\alpha)^2(3\alpha^2-2\alpha+2x-1)(4\alpha^2-\alpha+4x-2)}{4\alpha}$,
\[
\frac{\partial \Phi_J}{\partial \alpha}
=\frac{(\alpha-1)}{2\alpha^2}N_J(\alpha,x),
\quad
\begin{aligned}
N_J={}&30\alpha^5-40\alpha^4+30\alpha^3x-\alpha^3
-20\alpha^2x+9\alpha^2\\
&+4\alpha x^2-4\alpha x+\alpha+4x^2-4x+1.
\end{aligned}
\]
Substituting $x=\frac12+\alpha-\frac32\alpha^2+\alpha^2q$ into $N_P,N_J$ gives the
stationarity equations of \S\ref{ssec:branches}.

\section{Companion verification scripts}
\label{app:scripts}

All scripts live in \texttt{scripts/} and are run with
\texttt{/opt/miniconda3/bin/python3} (which has \texttt{sympy} 1.14; the system
\texttt{/usr/bin/python3} does \emph{not}).  Each uses exact rational arithmetic;
floating point appears only as a sanity print.

\begin{center}
\begin{tabular}{l|l}
script & certifies \\ \hline
\texttt{verify\_reduced\_problems.py} & branch uniqueness, monotonicity, endpoints, curvature (\S\ref{sec:certificates}) \\
\texttt{verify\_gap.py} & stationary-vs-boundary gap, Theorem~\ref{thm:stat-vs-bdry} \\
\texttt{verify\_balancing\_reduction.py} & Step~A2: densities, counterexample, saddle, Mantel (\S\ref{sec:balancing}) \\
\texttt{derive\_unbalanced.py} & SBM cross-check of the unbalanced $t(P),t(J)$ \\
\texttt{explore\_global\_infamily.py} & numpy scan: global in-family min $=$ balanced min \\
\texttt{optimize\_graphon.py} & direct step-graphon optimiser, Step~A1 evidence (\S\ref{sec:numerics}) \\
\end{tabular}
\end{center}

% ============================================================================
\begin{thebibliography}{9}

\bibitem{ReiherCliqueDensity}
C. Reiher,
\emph{The clique density theorem},
Annals of Mathematics \textbf{184} (2016), 683--707.
Also arXiv:1212.2454.

\bibitem{BennettDudekLidickyPikhurkoC5}
P. Bennett, A. Dudek, B. Lidick\'y, and O. Pikhurko,
\emph{Minimizing the number of 5-cycles in graphs with given edge-density},
Combinatorics, Probability and Computing \textbf{29} (2020), 44--67.
Also arXiv:1803.00165.

\bibitem{PikhurkoRazborov}
O. Pikhurko and A. Razborov,
\emph{Asymptotic structure of graphs with the minimum number of triangles},
Combinatorics, Probability and Computing \textbf{26} (2017), 138--160.
Also arXiv:1204.2846.

\bibitem{LovaszBook}
L. Lov\'asz,
\emph{Large Networks and Graph Limits},
American Mathematical Society, 2012.

\end{thebibliography}

\end{document}
